% ==============================================================================
% clean_federated_overleaf.tex
%
% Clean, compilable thesis starter file. Currently contains only the
% Centralised Baseline section. Build up the thesis by adding additional
% \section{...} blocks within the Results chapter as each experiment is
% written.
%
% Reporting style: descriptive (mean +/- SD across seeds), bolded best
% where comparing methods, no formal hypothesis testing -- matches the FL
% literature convention.
%
% To add a new section in the future:
%   1. Insert \section{Your section title} after the existing one.
%   2. Add \subsection{...} blocks for setup, results, role, summary.
%   3. Place tables inline as in this template, or split into a tables/
%      subfolder and use \input{tables/T01_xxx}.
%   4. Cross-references will resolve once corresponding \label{} markers
%      are added in other sections.
%
% Compiles standalone with pdflatex; no external .bib required yet.
% Bibliography will be added later when references are finalised.
% ==============================================================================

\documentclass[11pt,a4paper]{report}

\usepackage[utf8]{inputenc}
\usepackage[T1]{fontenc}
\usepackage{graphicx}
\usepackage{booktabs}
\usepackage{amsmath, amssymb}
\usepackage{bm}
\usepackage{microtype}
\usepackage{hyperref}
\usepackage{caption}
\usepackage{xcolor}
\usepackage[margin=2.5cm]{geometry}
\usepackage{setspace}
\onehalfspacing

\title{Federated Optimisation on DermaMNIST:\\
A Paired-Seed Empirical Comparison of FedAvg and FedProx\\
Under Statistical and Computational-Work Heterogeneity}
\author{Barbara Koch}
\date{}

\begin{document}

\maketitle

\tableofcontents
\newpage

% ==============================================================================
\chapter{Methods}
\label{ch:methods}
% ==============================================================================

This chapter describes the dataset, model, federated-learning
setup, algorithms, evaluation protocol, and implementation used
throughout the dissertation. The descriptions are at the
protocol-and-algorithm level; precise hyperparameter values are
collected in summary tables. The training/evaluation pipeline is
held identical across all experiments unless explicitly noted in a
given results section.

% ==============================================================================
\section{Dataset}
\label{sec:methods-dataset}
% ==============================================================================

All experiments use DermaMNIST, the dermatoscopic-imaging
benchmark from the MedMNIST~v2 collection. DermaMNIST is a
seven-class classification task on dermatoscopic skin-lesion
images sampled from HAM10000; representative samples from each
class are shown in Figure~\ref{fig:dataset-classes}. The
benchmark provides a canonical training/validation/test split,
summarised in Table~\ref{tab:dataset-splits}.

\begin{figure}[!htbp]
\centering
\includegraphics[width=\textwidth]{F_dataset_classes.pdf}
\caption{Representative training sample from each of the seven
DermaMNIST classes.}
\label{fig:dataset-classes}
\end{figure}

\begin{table}[!htbp]
\centering
\caption{DermaMNIST split sizes used throughout the dissertation.}
\label{tab:dataset-splits}
\small
\begin{tabular}{lc}
\toprule
Split & Samples \\
\midrule
Training   & $7{,}007$ \\
Validation & $1{,}003$ \\
Test       & $2{,}005$ \\
\midrule
Total      & $10{,}015$ \\
\bottomrule
\end{tabular}
\end{table}

\paragraph{Class distribution.}
The seven DermaMNIST classes are listed in
Table~\ref{tab:dataset-classes} together with their training-set
prevalence. The distribution is heavily skewed: the majority class
\emph{melanocytic nevi} accounts for $67.05\%$ of the training
set, while the rarest class \emph{dermatofibroma} accounts for
only $1.15\%$. This class imbalance is a feature of clinical
dermatology benchmarks (it reflects real-world incidence rates)
and is preserved without resampling so that all reported metrics
reflect performance on a realistic prior.

\begin{table}[!htbp]
\centering
\caption{DermaMNIST classes and training-set prevalence.}
\label{tab:dataset-classes}
\small
\begin{tabular}{lcc}
\toprule
Class & Training samples & Prevalence \\
\midrule
Actinic keratoses          & $228$   & $3.27\%$  \\
Basal cell carcinoma       & $359$   & $5.13\%$  \\
Benign keratosis           & $769$   & $10.97\%$ \\
Dermatofibroma             & $80$    & $1.15\%$  \\
Melanoma                   & $779$   & $11.11\%$ \\
Melanocytic nevi           & $4{,}693$ & $67.05\%$ \\
Vascular lesions           & $99$    & $1.41\%$  \\
\bottomrule
\end{tabular}
\end{table}

\paragraph{Preprocessing.}
DermaMNIST is loaded from the canonical \texttt{dermamnist\_64.npz}
distribution at $64\times64$ RGB resolution and resized to
$28\times28$ for all experiments. The downsampling is a
compute-economy choice: at $28\times28$ each communication round
completes in $\sim 3$--$5$~seconds on an A100, making ten paired
seeds tractable on a single-GPU SLURM allocation. The $28\times28$
resolution is the historical MedMNIST~v1 default and matches the
input resolution at which prior FL-on-DermaMNIST results in the
literature are reported. Pixel values are normalised using the
canonical ImageNet mean and standard deviation.

The $28\times28$ resolution heavily compresses dermatology-relevant
visual features such as pigment networks and asymmetry boundaries,
so absolute test performance reported in this dissertation should
be read as benchmark performance for FL-algorithm comparison rather
than as a clinical-grade dermatology estimate.

% ==============================================================================
\section{Model Architecture}
\label{sec:methods-model}
% ==============================================================================

The classifier used throughout this dissertation is
\texttt{DermMNISTCNN}, a four-block convolutional neural network
with $423{,}175$ trainable parameters. Each block consists of one
$3\times3$ convolution (\texttt{padding=1}) followed by a
GroupNorm normalisation layer, a ReLU activation, and a
$2\times2$ max-pooling operation; the final block replaces
max-pooling with adaptive global average pooling. The four
convolutional blocks expand the feature dimension as
$3 \to 32 \to 64 \to 128 \to 256$; a small fully-connected head
($256 \to 128 \to 7$) with dropout $0.2$ produces the class
logits. Table~\ref{tab:model-summary} summarises the architecture.

\begin{table}[!htbp]
\centering
\caption{\texttt{DermMNISTCNN} architecture summary. Each
convolutional block applies a $3\times3$ convolution (padding
$1$, stride $1$), GroupNorm with group count $g$, ReLU, and
spatial pooling. Input: $3 \times 28 \times 28$ RGB image;
output: $7$-class logits.}
\label{tab:model-summary}
\small
\begin{tabular}{llc}
\toprule
Stage & Operations & Output \\
\midrule
Block 1 & Conv $3 \to 32$ $+$ GN $g{=}4$ $+$ ReLU $+$ MaxPool $2\times2$    & $32 \times 14 \times 14$ \\
Block 2 & Conv $32 \to 64$ $+$ GN $g{=}8$ $+$ ReLU $+$ MaxPool $2\times2$   & $64 \times 7 \times 7$ \\
Block 3 & Conv $64 \to 128$ $+$ GN $g{=}16$ $+$ ReLU $+$ MaxPool $2\times2$ & $128 \times 3 \times 3$ \\
Block 4 & Conv $128 \to 256$ $+$ GN $g{=}16$ $+$ ReLU $+$ AdaptiveAvgPool   & $256$ \\
Head    & Linear $256 \to 128$ $+$ ReLU $+$ Dropout $0.2$ $+$ Linear $128 \to 7$ & $7$ \\
\bottomrule
\end{tabular}
\end{table}

\paragraph{GroupNorm rather than BatchNorm.}
The choice of GroupNorm rather than BatchNorm is deliberate.
Per-client BatchNorm statistics encode that client's local data
distribution; when the global model is constructed by parameter
averaging, these batch statistics are averaged across heterogeneous
clients in a way that can corrupt the running estimates and harm
generalisation. Hsieh et al.\ (2020) document this effect and
recommend normalisation layers whose statistics are computed
within each example (GroupNorm, LayerNorm) for federated learning
on non-IID data. We adopt GroupNorm following this recommendation.

% ==============================================================================
\section{Federated Learning Setup}
\label{sec:methods-fl-setup}
% ==============================================================================

The federated learning protocol is cross-silo and synchronous,
matching the canonical formulation surveyed in Kairouz et al.\
(2021). The training set is partitioned across $K = 7$ simulated
clients (the per-experiment partition assignment is described
in the relevant Results section). At each communication round
$t$, the central server broadcasts the current global model
parameters $w^t$ to all clients; each participating client $i$
performs $E$ epochs of local SGD on its local training subset
$\mathcal{D}_i$, yielding updated parameters $w_i^{t+1}$. The
server then forms the next global iterate by size-weighted
parameter averaging:
\[
w^{t+1} = \sum_{i=1}^{K} \frac{n_i}{\sum_j n_j}\, w_i^{t+1},
\]
where $n_i = |\mathcal{D}_i|$ is the number of training samples
held by client $i$. The size-weighting rule is the standard
FedAvg aggregation of McMahan et al.\ (2017). Table~\ref{tab:fl-setup}
collects the federated-learning hyperparameters.

\begin{table}[!htbp]
\centering
\caption{Federated learning hyperparameters held constant across
experiments unless explicitly noted in the relevant Results
section.}
\label{tab:fl-setup}
\small
\begin{tabular}{lc}
\toprule
Hyperparameter & Value \\
\midrule
Number of clients $K$                & $7$ \\
Client participation $C$             & $1.0$ (all clients per round) \\
Communication rounds $R$             & $150$ \\
Local epochs $E$                     & $20$ \\
Aggregation rule                     & Size-weighted parameter averaging \\
\bottomrule
\end{tabular}
\end{table}

\paragraph{Client count.}
The choice of $K = 7$ clients reflects a deliberate trade-off
between simulating multi-institutional federation and keeping
each client's training set large enough to support stable local
optimisation. With $K = 7$ and $7{,}007$ training samples, each
client receives on the order of $10^3$ samples, which is
sufficient for the 423K-parameter \texttt{DermMNISTCNN} to learn
stably; with $K \gg 10$ the per-client sample count would fall
below the regime in which an architecture of this capacity
converges reliably at $28\times28$ resolution.

% ==============================================================================
\section{Algorithms}
\label{sec:methods-algorithms}
% ==============================================================================

Two federated optimisation algorithms are compared throughout
the dissertation: FedAvg (McMahan et al., 2017) and FedProx
(Li et al., 2020). Both use the same server-side size-weighted
parameter averaging defined above; they differ only in the local
training objective each client minimises.

\paragraph{FedAvg.}
Each client $i$ minimises the standard local cross-entropy loss
on its local data:
\[
\min_{w} \; F_i(w) \;=\; \frac{1}{n_i} \sum_{(x, y) \in \mathcal{D}_i}
\ell\big(f_w(x),\, y\big),
\]
where $\ell$ is the cross-entropy loss and $f_w$ is the network
parameterised by $w$. Local optimisation uses SGD for $E$ epochs
starting from the broadcast global parameters $w^t$.

\paragraph{FedProx.}
Each client $i$ minimises an augmented local objective that adds
a proximal regulariser pulling the local parameters toward the
broadcast global anchor:
\[
\min_{w} \; F_i(w) \;+\; \frac{\mu}{2} \lVert w - w^t \rVert_2^2,
\]
where $\mu \geq 0$ is the proximal coefficient. The first term is
the local objective minimised by FedAvg; the second is the
proximal regulariser. Setting $\mu = 0$ recovers the FedAvg local
objective exactly, so the two algorithms differ only through the
proximal term when $\mu > 0$. Throughout the main experiments we
set $\mu = 0.01$, the value reported by Li et al.\ (2020) as a
robust default across their benchmark matrix; the choice is
fixed a priori and is not tuned per experiment.

\paragraph{Equivalence check at $\mu = 0$.}
The FedProx implementation contains an explicit zero-mu branch
that bypasses the proximal gradient contribution when $\mu = 0$,
making the FedProx local-training loop bit-equivalent to the
FedAvg loop in that regime. This equivalence is verified by
running both algorithms with $\mu = 0$ on the same seed and
confirming identical parameter trajectories. The two algorithms
therefore differ \emph{only} through the $\mu > 0$ proximal
contribution.

% ==============================================================================
\section{Training and Evaluation Protocol}
\label{sec:methods-protocol}
% ==============================================================================

\paragraph{Optimiser and local training.}
All clients use stochastic gradient descent with the
hyperparameters in Table~\ref{tab:training-hparams}. The
local-training loop is identical across both algorithms: $E$
epochs of mini-batched SGD on the client's local training
subset, with the local loss function differing only through the
optional proximal term (Section~\ref{sec:methods-algorithms}).

\begin{table}[!htbp]
\centering
\caption{Optimiser and local-training hyperparameters.}
\label{tab:training-hparams}
\small
\begin{tabular}{lc}
\toprule
Hyperparameter & Value \\
\midrule
Optimiser                        & SGD \\
Learning rate                    & $0.01$ \\
Momentum                         & $0.9$ \\
Weight decay                     & $0$ \\
Batch size                       & $32$ \\
Local epochs per round $E$       & $20$ \\
Local loss                       & Cross-entropy \\
FedProx proximal coefficient $\mu$ & $0.01$ \\
\bottomrule
\end{tabular}
\end{table}

\paragraph{Evaluation.}
The server evaluates the global model on the centralised
validation set after every communication round, recording
validation cross-entropy loss, validation accuracy, validation
balanced accuracy, validation macro-F1, and per-class
validation F1. At the end of training, the test set is
evaluated \emph{once} on the global model produced at the
round of highest validation macro-F1 (the
\emph{best-validation checkpoint} protocol). Test-set
macro-F1 at this checkpoint is the headline metric reported
throughout the dissertation; per-class F1, balanced accuracy,
and overall accuracy on the test set are reported as
secondary measures. The validation set is held fixed across
all experiments; the test set is never seen by the model
during training and is evaluated exactly once per run.

% ==============================================================================
\section{Paired-Seed Protocol}
\label{sec:methods-paired-seed}
% ==============================================================================

A defining design choice of the experimental protocol used in
this dissertation is the \emph{paired-seed} comparison
structure. For every FedAvg-vs-FedProx contrast reported in
the Results chapter, the two algorithms are run on the same
set of ten random seeds: $\{42,\, 123,\, 456,\, 789,\, 999,\,
2024,\, 31337,\, 8675309,\, 161803,\, 271828\}$. Within each
seed, the following sources of stochasticity are fixed to
identical values across the FedAvg and FedProx runs:

\begin{itemize}
\setlength\itemsep{2pt}
\item the random partition of training samples to clients;
\item the initial model parameters $w^0$;
\item the per-(round, client) dataloader shuffle order;
\item the per-(round, client) local-epoch schedule, where
      applicable (for the system-heterogeneity experiments
      in Chapter~\ref{ch:results});
\item the random sampling of participating clients per round,
      where applicable.
\end{itemize}

Because all these sources of stochasticity are matched
within-pair, the within-pair difference
\[
\Delta_s \;=\; \text{test macro-F1}_{\text{FedProx}, s}
              - \text{test macro-F1}_{\text{FedAvg}, s}
\]
on seed $s$ reflects only the algorithmic difference (the
gated proximal term), not random-state divergence. Aggregate
within-pair differences across the ten seeds therefore have
substantially tighter variance than the corresponding
independent-seed comparison would, allowing meaningful
estimates of small algorithmic effects at modest sample size.

\paragraph{Reporting style.}
We follow the convention of the federated-learning literature
(McMahan et al.\ 2017, Li et al.\ 2020, Wang et al.\ 2020,
Hsu et al.\ 2019, Acar et al.\ 2021) in reporting paired
results descriptively: mean $\pm$ standard deviation across
seeds, visual curve comparisons with shaded $\pm$~SEM bands,
bolded best results in tables, and within-pair $\Delta$
values with the count of paired-seed wins. Formal hypothesis
testing is not used in the Results chapter; differences
smaller in magnitude than the across-seed standard deviation
are described as ``descriptively positive/negative'' rather
than as significant or null effects.

% ==============================================================================
\section{Implementation and Compute}
\label{sec:methods-implementation}
% ==============================================================================

The experiments are implemented in two complementary code paths
that share the same algorithmic logic but use different runtimes.

\paragraph{Pure-PyTorch reference loop.}
A single-process sequential implementation in PyTorch that
iterates through clients in deterministic order, holds the
global parameter state in-memory, and performs explicit
size-weighted aggregation after each round. This implementation
is used as the algorithmic reference and is the primary runtime
for the headline statistical-heterogeneity result
(Section~\ref{sec:engineered-headline}).

\paragraph{Flower simulation runtime.}
A federated-learning simulation using the Flower framework
(Beutel et al., 2020) with each client implemented as a Ray
actor. The server-side aggregation passes through Flower's
\texttt{FedAvg} strategy with size-weighted averaging; the
client-side FedAvg/FedProx local training loop is functionally
identical to the pure-PyTorch path. This runtime is used for
the cross-runtime sensitivity check
(Section~\ref{sec:runtime-sensitivity}) and for all
robustness, system-heterogeneity, and partial-participation
experiments in the Results chapter.

The two implementations have been verified to use bit-equivalent
local objectives at $\mu = 0$ (the FedProx code path collapses
to FedAvg in the $\mu = 0$ branch in both runtimes). Differences
in per-seed results between the two runtimes therefore arise
from implementation-level factors --- floating-point
non-associativity of summation order in aggregation, parameter
serialisation across Ray actor process boundaries, and
random-number-generator routing through Ray scheduling ---
rather than from algorithmic differences. The cross-runtime
sensitivity check in Section~\ref{sec:runtime-sensitivity}
characterises the magnitude of this implementation-level
variance.

\paragraph{Compute.}
Pure-PyTorch experiments are executed on the Cambridge HPC
ampere partition (NVIDIA A100~80GB). Flower-runtime
experiments are executed on either the same Cambridge HPC
partition or on rented A100/H100 instances on RunPod where
queue priority constrains turnaround. All result JSON files
record the host machine, CUDA device, framework version, git
commit, and start/finish timestamps for full reproducibility.
A single training run completes in approximately $5$~minutes
(centralised) to $20$~minutes (federated $150$-round) on an
A100, with FedProx runs taking $\sim 40\%$ longer than FedAvg
runs at matched conditions due to the per-step computation of
the proximal regulariser.

\paragraph{Reproducibility.}
The complete experimental matrix can be reproduced from the
SLURM submission scripts in \texttt{mnist\_dermnist/scripts/}
in the project repository. Each script declares the partition,
seed list, algorithm, and hyperparameters explicitly; together
with the run-time provenance recorded in each result JSON,
this makes every reported number traceable to a specific
submission script, git commit, and machine.

% ==============================================================================
\chapter{Results and Discussion}
\label{ch:results}
% ==============================================================================

This chapter reports the empirical findings of the dissertation,
starting with a centralised non-federated reference and then
progressing through the federated experiments. Methodological
observations and mechanism interpretation are interleaved with the
empirical results rather than relegated to a separate discussion
chapter.

% ==============================================================================
\section{Centralised Baseline}
\label{sec:centralised-baseline}
% ==============================================================================

Before evaluating any federated configuration, we establish a
centralised non-federated reference. The centralised baseline trains
the same \texttt{DermMNISTCNN} architecture used in the federated
experiments on the pooled DermaMNIST training set with no
federation. It provides an upper bound on what can be achieved on
this task with this architecture, against which the federated
results reported in the remainder of this chapter are interpreted
as fractional recoveries of the centralised performance.

\subsection{Training setup}
\label{sec:centralised-setup}

The centralised baseline pools all $7{,}007$ training samples,
$1{,}003$ validation samples, and $2{,}005$ test samples of the
DermaMNIST benchmark into a single non-federated training loop.
The model architecture is held identical to the federated
experiments: a four-block GroupNorm CNN with $423{,}175$ trainable
parameters operating on $28\times28$ RGB inputs. The optimiser is
SGD with learning rate $0.01$, momentum $0.9$, weight decay $0$,
and batch size $32$ --- the same hyperparameters used by each
client in the subsequent federated experiments, chosen for direct
comparability rather than for centralised optimum. The model is
trained for $50$ epochs; test set performance is reported at the
epoch of highest validation macro-F1 (the
\emph{best-validation checkpoint} protocol).

To match the federated experiments and enable paired-seed analyses
where applicable, the centralised baseline is repeated across the
same ten random seeds used for all federated runs:
$\{42,\,123,\,456,\,789,\,999,\,2024,\,31337,\,8675309,\,161803,\,$
$271828\}$. Within each seed the data shuffle, model
initialisation, and training trajectory are deterministic; the
across-seed variation quantifies optimisation noise on this task.

\subsection{Results}
\label{sec:centralised-results}

Table~\ref{tab:centralised-results} reports test performance at
the best-validation checkpoint across the ten seeds. The
centralised model reaches a mean test macro-F1 of
$0.560\pm0.024$, mean classification accuracy of
$0.764\pm0.012$, and mean balanced accuracy of
$0.561\pm0.034$. The per-seed test macro-F1 ranges from
$0.518$ to $0.597$. The best-validation epoch is consistently in
the second half of training (median $47$ out of $50$ epochs),
indicating that the $50$-epoch training budget is appropriate
for this task and architecture.

\begin{table}[!htbp]
\centering
\caption{Centralised baseline test performance at the
best-validation checkpoint, reported as mean $\pm$ standard
deviation across $n=10$ seeds. Per-class F1 is the F1 score on
each of the seven DermaMNIST classes; class prevalence in the
training set is shown in parentheses.}
\label{tab:centralised-results}
\small
\begin{tabular}{lc}
\toprule
Metric & Test value (mean $\pm$ SD) \\
\midrule
Macro-F1                        & $\mathbf{0.5600 \pm 0.0240}$ \\
Balanced accuracy               & $0.5608 \pm 0.0341$ \\
Overall accuracy                & $0.7637 \pm 0.0123$ \\
Cross-entropy loss              & $0.7778 \pm 0.0653$ \\
\midrule
\multicolumn{2}{l}{\textit{Per-class F1:}} \\
\quad Actinic keratoses (3.27\%)              & $0.4923 \pm 0.0433$ \\
\quad Basal cell carcinoma (5.13\%)           & $0.5380 \pm 0.0263$ \\
\quad Benign keratosis (10.97\%)              & $0.5285 \pm 0.0397$ \\
\quad Dermatofibroma (1.15\%)                 & $0.3359 \pm 0.0961$ \\
\quad Melanoma (11.11\%)                      & $0.4551 \pm 0.0349$ \\
\quad Melanocytic nevi (67.05\%)              & $0.8896 \pm 0.0088$ \\
\quad Vascular lesions (1.41\%)               & $0.6807 \pm 0.0680$ \\
\bottomrule
\end{tabular}
\end{table}

The per-class breakdown reveals the task's structure clearly.
The majority class \emph{melanocytic nevi} ($67\%$ prevalence)
is the easiest to learn, reaching $F1 = 0.890$ with very small
across-seed variance ($\pm 0.009$). The two rarest classes,
\emph{dermatofibroma} ($1.15\%$ prevalence, $80$ training
samples) and \emph{vascular lesions} ($1.41\%$ prevalence,
$99$ training samples), show the largest across-seed variance
($\pm 0.096$ and $\pm 0.068$), reflecting the difficulty of
stably learning classes represented by such small absolute
counts. The two clinically critical classes \emph{melanoma}
and \emph{basal cell carcinoma} reach mean F1 of $0.455$ and
$0.538$ respectively --- well above chance but substantially
below the majority class, consistent with the medical-imaging
literature on small CNNs at $28\times28$ resolution.

The macro-F1 of $0.56$ is a benchmark-typical result for
DermaMNIST at this resolution and architecture, comparable to
published values on the MedMNIST v2 benchmark. It is not
intended as a clinical-grade performance estimate; the input
resolution is deliberately compressed and the architecture is a
small CNN, both chosen for comparability across the federated
experiments in this dissertation rather than for clinical
deployment.

\subsection{Role in the subsequent analysis}
\label{sec:centralised-role}

The centralised baseline serves three purposes throughout the
remainder of this chapter.

First, it is the \emph{upper bound} against which federated
configurations are measured. Subsequent results report recovery
percentages relative to the $0.5600$ centralised mean
(e.g.\ ``FedProx recovers $90.7\%$ of centralised performance''),
which makes the absolute magnitude of federation costs
interpretable. Within-federation algorithmic differences are
then framed against this fixed reference rather than against
each other in isolation.

Second, the across-seed standard deviation of the centralised
baseline ($\pm 0.024$ macro-F1) sets a noise floor below which
seed-level differences in federated experiments cannot be
reliably distinguished from optimisation noise alone. Federated
within-pair differences smaller than this in magnitude are read
as descriptive rather than substantive in the remainder of this
chapter.

Third, the per-class profile in
Table~\ref{tab:centralised-results} establishes the
\emph{intrinsic difficulty} of each class on this task,
independent of any federation effect. Federated per-class
breakdowns reported later are interpreted relative to this
profile: a class that already shows high variance in the
centralised setting is unlikely to support strong inferential
claims in the federated setting either.

\subsection{Summary}
\label{sec:centralised-summary}

The centralised baseline establishes that the $28\times28$
DermaMNIST classification task with the \texttt{DermMNISTCNN}
architecture supports a mean test macro-F1 of approximately
$0.56$ under non-federated training, with the majority class
melanocytic nevi at near-ceiling performance ($F1 \approx 0.89$)
and the rarest classes at substantially lower and noisier
per-class F1. This serves as the reference against which all
federated results in the rest of this chapter are interpreted.

% ==============================================================================
\section{IID Federated Baseline}
\label{sec:iid-baseline}
% ==============================================================================

Before evaluating the FedProx-vs-FedAvg comparison under
heterogeneous client data distributions, we establish a federated
baseline in which every client holds an i.i.d.\ sample of the
global training distribution. The IID configuration is the
\emph{mechanism null} condition: when local objectives $F_i$ are
i.i.d.\ samples of the same underlying distribution, the bounded-
dissimilarity argument of Li et al.\ (2020)~\S4.1 predicts that
no inter-client drift exists for the proximal anchor to constrain
and FedProx should provide no benefit over FedAvg. A positive
$\Delta$ under IID would therefore contradict the mechanism
interpretation that the subsequent non-IID experiments rely on.
This section confirms the mechanism null empirically.

\subsection{Partition and training setup}
\label{sec:iid-setup}

The IID partition (\texttt{iid\_7\_clients}) assigns the
$7{,}007$ training samples to seven simulated clients by uniform
random shuffling of the global index list and splitting into
seven equal shards. Each client therefore receives the same
number of samples and, in expectation, the same class
composition as the global training set --- so client-local
distributions differ from one another only through finite-sample
sampling noise, not by design.

Figure~\ref{fig:iid-partition} visualises the resulting
per-(client, class) sample counts at seed $42$. All seven
clients hold exactly $1{,}001$ samples, and each client's
stacked composition mirrors the global class distribution within
the small variation expected from random partitioning of a
$7{,}007$-sample pool. In absolute counts, every client receives
approximately $670$ melanocytic nevi, $110$ melanoma, $110$
benign keratosis, $50$ basal cell carcinoma, $33$ actinic
keratoses, $14$ vascular lesions, and $11$ dermatofibroma
samples; all seven classes are present in every client, including
the three rarest ones. The IID partition therefore represents
the limit case in which no client carries class-specialised
information: any class that a client must learn to predict, it
already sees locally in roughly the global proportion.

\begin{figure}[!htbp]
\centering
\includegraphics[width=0.95\textwidth]{F_iid_partition.pdf}
\caption{IID partition of the $7{,}007$-sample DermaMNIST
training set across the seven simulated clients (seed $42$).
Each bar is one client; segment heights show the number of
samples of each class. All clients receive exactly $1{,}001$
samples and the stacks are visually near-identical, confirming
that the per-client class composition matches the global
distribution up to finite-sample sampling noise.}
\label{fig:iid-partition}
\end{figure}

All other hyperparameters match the centralised baseline of
Section~\ref{sec:centralised-baseline}: identical
\texttt{DermMNISTCNN} architecture, SGD with learning rate
$0.01$, momentum $0.9$, batch size $32$, weight decay $0$.
Federated-specific hyperparameters are: full participation
($C = 1.0$, all seven clients per round), $R = 150$ communication
rounds, $E = 20$ local epochs, FedProx proximal coefficient
$\mu = 0.01$. The Flower simulation framework is used as the
canonical FL runtime throughout. The same ten paired seeds as
the centralised baseline are used; within each seed the partition,
initial model weights, and client-side data shuffling are
deterministic and identical across the FedAvg and FedProx arms,
so within-pair differences reflect only the algorithmic gating
of the proximal term.

\subsection{Results}
\label{sec:iid-results}

Table~\ref{tab:iid-results} reports test performance at the
best-validation checkpoint across the ten seeds. FedAvg and
FedProx reach indistinguishable mean test macro-F1
($0.585 \pm 0.020$ and $0.579 \pm 0.019$ respectively); the
within-pair mean difference is $-0.006 \pm 0.019$ and FedProx
wins on $5$ of $10$ paired seeds.

\begin{table}[!htbp]
\centering
\caption{IID federated baseline. Test performance at the
best-validation checkpoint, mean $\pm$ standard deviation across
$n = 10$ paired seeds. The winning algorithm per row is shown in
bold; for rows where the two algorithms are within $\pm 0.001$
neither is bolded. Per-class F1 is the F1 score on each of the
seven DermaMNIST classes, with class prevalence in the training
set in parentheses --- these prevalences are global and, under
IID partitioning, also describe each client's expected local
class distribution (see Figure~\ref{fig:iid-partition}).}
\label{tab:iid-results}
\small
\begin{tabular}{lccc}
\toprule
Metric & FedAvg & FedProx & $\Delta$ \\
\midrule
Test macro-F1                 & $\mathbf{0.5846 \pm 0.0197}$ & $0.5786 \pm 0.0190$ & $-0.0061$ \\
Balanced accuracy             & $\mathbf{0.5852 \pm 0.0290}$ & $0.5802 \pm 0.0245$ & $-0.0050$ \\
Overall accuracy              & $0.7625 \pm 0.0094$          & $0.7613 \pm 0.0083$ & $-0.0012$ \\
\midrule
\multicolumn{4}{l}{\textit{Per-class F1:}} \\
\quad Actinic keratoses (3.27\%)   & $\mathbf{0.523 \pm 0.034}$ & $0.507 \pm 0.030$ & $-0.016$ \\
\quad Basal cell carcinoma (5.13\%)& $\mathbf{0.557 \pm 0.011}$ & $0.542 \pm 0.034$ & $-0.015$ \\
\quad Benign keratosis (10.97\%)   & $0.538 \pm 0.024$          & $0.537 \pm 0.020$ & $-0.001$ \\
\quad Dermatofibroma (1.15\%)      & $\mathbf{0.412 \pm 0.075}$ & $0.406 \pm 0.069$ & $-0.006$ \\
\quad Melanoma (11.11\%)           & $\mathbf{0.462 \pm 0.020}$ & $0.456 \pm 0.018$ & $-0.007$ \\
\quad Melanocytic nevi (67.05\%)   & $0.894 \pm 0.004$          & $0.894 \pm 0.004$ & $+0.000$ \\
\quad Vascular lesions (1.41\%)    & $0.707 \pm 0.040$          & $0.708 \pm 0.036$ & $+0.001$ \\
\bottomrule
\end{tabular}
\end{table}

The two algorithms are descriptively indistinguishable across
all reported metrics. The mean within-pair difference is
small and negative ($-0.006$ macro-F1, $-0.005$ balanced
accuracy) but lies within the across-seed standard deviation of
either algorithm. The per-class breakdown shows the same pattern
on every class: differences are at most $0.016$ in absolute
value and have no consistent direction. Notably, the majority
class \emph{melanocytic nevi} produces a per-pair difference
of $+0.0001$ macro-F1 with FedProx and FedAvg matching to three
decimal places ($0.894$ vs $0.894$); this is the most
mechanistically transparent class because it is structurally
held by every client at the same proportion, and its near-zero
$\Delta$ is the expected signature of an inert proximal anchor.

Figure~\ref{fig:iid-convergence} shows the per-round validation
macro-F1 trajectories on the IID partition; the FedAvg and FedProx
curves are essentially superimposed across all 150 rounds,
visually confirming the null $\Delta$ reported in
Table~\ref{tab:iid-results}.

\begin{figure}[!htbp]
\centering
\includegraphics[width=0.85\textwidth]{F_iid_convergence.pdf}
\caption{Validation macro-F1 over 150 communication rounds on
the IID partition (Flower runtime, $n = 10$ paired seeds,
mean $\pm$ SEM band). FedAvg and FedProx produce essentially
identical trajectories from round $\sim 10$ onwards, consistent
with the proximal-anchor mechanism's prediction that the
regulariser is inert when no inter-client class skew exists.}
\label{fig:iid-convergence}
\end{figure}

An independent IID replication is available in
\texttt{results/iid/}, which uses the same partition,
hyperparameters, and ten seeds but was run as a separate sweep.
The independent replication produces a mean within-pair
difference of $-0.007 \pm 0.022$ macro-F1 with FedProx winning
$4$ of $10$ paired seeds, agreeing with the canonical IID
baseline to within $0.001$ macro-F1. The two independent sweeps
converging on the same near-zero $\Delta$ provides additional
evidence that the IID null is genuine rather than a feature of
any single sweep's random seed draws.

\subsection{Role in the subsequent analysis}
\label{sec:iid-role}

The IID baseline serves two roles in the remainder of the
chapter. First, it provides the empirical confirmation of the
\emph{mechanism null prediction}: in a setting where local
objectives are i.i.d.\ samples of the same global distribution,
the proximal anchor has no inter-client drift to constrain and
FedProx does not differ meaningfully from FedAvg. Any positive
within-pair $\Delta$ observed in subsequent non-IID experiments
can therefore be interpreted as evidence of mechanism activation
relative to this null, rather than as an artefact of the
algorithm or the architecture.

Second, the IID across-seed standard deviation of approximately
$\pm 0.020$ macro-F1 sets a noise floor for FL paired-comparison
analyses on this task at $n = 10$ seeds. Within-pair $\Delta$
values of comparable or smaller magnitude in subsequent
non-IID experiments are read as descriptive rather than as
substantive findings.

\subsection{Summary}
\label{sec:iid-summary}

Under uniform random partitioning of the DermaMNIST training set
into seven clients, FedAvg and FedProx produce
indistinguishable test performance (mean $\Delta = -0.006$
macro-F1, FedProx wins $5/10$ paired seeds; independent
replication: $\Delta = -0.007$, $4/10$). This is the predicted
behaviour of the proximal anchor mechanism when no inter-client
class skew exists. The IID baseline confirms that any positive
$\Delta$ observed in the subsequent non-IID experiments is
attributable to the partition design rather than to a
spurious algorithmic advantage of FedProx under any
partitioning.

% ==============================================================================
\section{Engineered Non-IID Partition: Primary Result}
\label{sec:engineered-headline}
% ==============================================================================

The IID baseline of Section~\ref{sec:iid-baseline} established
that without inter-client class skew, FedProx and FedAvg are
empirically indistinguishable. This section reports the primary
non-IID result of the dissertation: the FedAvg-vs-FedProx
comparison on the engineered \texttt{balanced\_paired\_7\_clients}
partition, which deliberately creates non-IID class skew
designed to expose the proximal anchor's drift-control
mechanism. Because cross-runtime sensitivity is a known
methodological concern in federated learning, the same paired
comparison is reported on two independent implementations of the
algorithms (a pure-PyTorch sequential reference loop and the
Flower simulation framework) so that runtime-induced variation
can be assessed alongside the headline effect.

\subsection{Partition design}
\label{sec:engineered-partition}

The engineered \texttt{balanced\_paired\_7\_clients} partition
assigns each of the six minority classes to exactly two of the
seven clients (a \emph{paired-co-training} design), with the
majority class \emph{melanocytic nevi} held in equal proportion
by all seven clients. The seven minority-class slots are paired
as follows: clients $C_0$ and $C_1$ jointly hold actinic
keratoses and basal cell carcinoma; $C_2$ and $C_3$ hold benign
keratosis and dermatofibroma; $C_4$ and $C_5$ hold melanoma and
vascular lesions; and $C_6$ is a nevi-only generalist that holds
no minority-class samples. Within each pair, samples of the
shared classes are split as evenly as the integer counts allow
(e.g.\ basal cell carcinoma is split $180$/$179$ between $C_0$
and $C_1$). Concretely, $C_0$ holds $114$ actinic keratoses
$+\,180$ basal cell carcinoma $+\,670$ mel-nevi ($n = 964$);
$C_2$ holds $385$ benign keratosis $+\,40$ dermatofibroma
$+\,670$ mel-nevi ($n = 1{,}095$); $C_4$ holds $390$ melanoma
$+\,50$ vascular lesions $+\,670$ mel-nevi ($n = 1{,}110$); and
$C_6$ holds $673$ mel-nevi only. Per-client totals therefore lie
between $673$ and $1{,}110$ samples --- a $1.65\times$ range
that keeps quantity skew intentionally mild relative to label
skew, so the partition isolates the label-skew axis as the
primary source of inter-client drift.

Figure~\ref{fig:engineered-partition} visualises the resulting
per-(client, class) sample counts in the same format as the IID
partition figure (Figure~\ref{fig:iid-partition}), enabling a
direct visual contrast: under IID every client's stack is a
scaled copy of the global distribution, whereas under the
engineered partition the stacks differ structurally --- four of
the seven clients hold zero samples of a given minority class.
Every minority class is therefore held by exactly two co-trained
clients that must drift apart during local training and be
reconciled at aggregation, while the majority class is held by
every client and provides a within-experiment per-class control
where the proximal anchor is mechanistically inert (cf.\ the
mel-nevi observation in the IID baseline).

\begin{figure}[!htbp]
\centering
\includegraphics[width=0.95\textwidth]{F_engineered_partition.pdf}
\caption{Engineered \texttt{balanced\_paired\_7\_clients}
partition of the $7{,}007$-sample DermaMNIST training set.
Each bar is one client; segment heights show the number of
samples of each class. Per-client totals range from $673$ to
$1{,}110$. Each minority class appears on exactly two clients
(its co-training pair), and four of the seven clients hold no
samples of any given minority class --- the structural label
skew that the FedProx proximal anchor is designed to control.
Direct visual comparison with Figure~\ref{fig:iid-partition}
makes the qualitative difference between the two partitions
explicit.}
\label{fig:engineered-partition}
\end{figure}

\subsection{Training setup}
\label{sec:engineered-setup}

All training hyperparameters are held identical to the IID
baseline of Section~\ref{sec:iid-setup}: SGD with learning rate
$0.01$, momentum $0.9$, batch size $32$, weight decay $0$,
$R = 150$ communication rounds, $E = 20$ local epochs, full
participation $C = 1.0$, FedProx proximal coefficient
$\mu = 0.01$, and the same ten paired seeds. The only variable
relative to the IID baseline is the partition design.

The paired comparison is executed under two independent
implementations of FedAvg and FedProx:

\begin{itemize}
\item \textbf{Pure-PyTorch reference loop:} a single-process
sequential implementation in \texttt{fl/server\_loop.py} with
deterministic client iteration, in-memory parameter passing, and
explicit size-weighted aggregation. This serves as the
algorithmic reference.
\item \textbf{Flower simulation framework:} the standard
\texttt{flwr.simulation} runtime with each client implemented
as a Ray actor. Parameter aggregation passes through Flower's
\texttt{FedAvg} strategy, with parameters serialised across
process boundaries.
\end{itemize}

Both implementations have been verified to use bit-equivalent
local objectives at $\mu = 0$ (the FedProx code path collapses
to FedAvg in the $\mu = 0$ branch in both runtimes) and to use
identical hyperparameters. Within each runtime, paired-seed
controls are applied so that within-pair differences reflect
only the algorithmic gating of the proximal term.

\subsection{Results}
\label{sec:engineered-results}

Table~\ref{tab:engineered-results} reports test performance at
the best-validation checkpoint across the ten paired seeds on
each runtime. Figure~\ref{fig:engineered-convergence} shows the
per-round validation macro-F1 trajectories.

\begin{table}[!htbp]
\centering
\caption{Engineered \texttt{balanced\_paired\_7\_clients}
partition. Test performance at the best-validation checkpoint,
mean $\pm$ standard deviation across $n = 10$ paired seeds. The
winning algorithm per runtime is shown in bold.}
\label{tab:engineered-results}
\small
\begin{tabular}{lccccc}
\toprule
Runtime & FedAvg & FedProx & $\Delta$ & FedProx wins \\
\midrule
Pure-PyTorch (reference) & $0.4814 \pm 0.0254$ & $\mathbf{0.5081 \pm 0.0140}$ & $+0.0267$ & $9/10$ \\
Flower (simulation)      & $0.4964 \pm 0.0150$ & $\mathbf{0.5033 \pm 0.0206}$ & $+0.0069$ & $7/10$ \\
\bottomrule
\end{tabular}
\end{table}

Under both runtimes the direction of effect is positive: FedProx
outperforms FedAvg in mean macro-F1, with FedProx winning the
majority of paired seeds. The magnitude differs substantially
between runtimes: the within-pair difference is $+0.027$ on
pure-PyTorch and $+0.007$ on Flower --- a factor of approximately
four. Per-seed wins are $9$ of $10$ under pure-PyTorch and $7$ of
$10$ under Flower; the lone losing seed under pure-PyTorch
($\Delta = -0.020$ on seed~789) shifts to a winning seed under
Flower, and three of the nine pure-PyTorch winning seeds shift
to losing seeds under Flower. The pure-PyTorch result also shows
a notably smaller across-seed standard deviation for FedProx
($\pm 0.014$) than for FedAvg ($\pm 0.025$); under Flower the SDs
are closer ($\pm 0.021$ vs $\pm 0.015$).

The convergence trajectories
(Figure~\ref{fig:engineered-convergence}) show that on both
runtimes FedProx maintains a higher validation macro-F1 than
FedAvg throughout training, with the gap visible from
round~$\sim 10$ onwards. The pure-PyTorch gap is consistently
larger across the entire training horizon, not just at the
best-validation checkpoint, consistent with the larger headline
$\Delta$ reported in Table~\ref{tab:engineered-results}.

\begin{figure}[!htbp]
\centering
\includegraphics[width=\textwidth]{F_engineered_convergence.pdf}
\caption{Validation macro-F1 over 150 communication rounds on
the engineered \texttt{balanced\_paired\_7\_clients} partition
under two independent implementations ($n = 10$ paired seeds,
mean $\pm$ SEM band). FedProx (dark slate blue) maintains a
higher trajectory than FedAvg (mint green) under both runtimes;
the gap is visibly larger under pure-PyTorch (left) than under
Flower (right).}
\label{fig:engineered-convergence}
\end{figure}

\subsection{Runtime sensitivity}
\label{sec:runtime-sensitivity}

The $4\times$ difference in headline $\Delta$ between runtimes
on the same partition with the same seeds is a substantively
larger gap than the within-runtime across-seed standard
deviations would predict for replicating implementations of the
same algorithm. To check whether this gap is a systematic
per-seed shift or a high-variance noise effect, we computed the
per-seed difference $\Delta_{\mathrm{PT},s} -
\Delta_{\mathrm{Flower},s}$. The per-seed differences are not
consistently signed: five seeds show $\Delta_{\mathrm{PT}} >
\Delta_{\mathrm{Flower}}$ and five show
$\Delta_{\mathrm{Flower}} \geq \Delta_{\mathrm{PT}}$. The
across-seed mean of the differences is $+0.020$, but the
distribution spans both signs --- indicating that the runtime
gap is a noisy mean-level effect rather than a systematic
per-seed shift.

Three implementation-level differences plausibly contribute:
the Flower runtime runs each client in a separate Ray actor
with parameter serialisation across process boundaries, which
introduces floating-point operations whose order is not
guaranteed identical to the single-process pure-PyTorch loop;
the per-(round, client) random-number-generator state is
threaded differently through Flower's scheduling than through
the pure-PyTorch loop, even when the same user-level seed is
supplied; and Flower's \texttt{FedAvg} aggregation strategy
sums client updates in an order determined by actor return
times rather than client identity, which makes the aggregated
parameter vector sensitive to floating-point non-associativity
across 150 rounds and 7 clients. None of these differences are
algorithmic --- the FedAvg/FedProx local objectives and the
size-weighted aggregation rule are bit-equivalent at the
mathematical level --- but they are sufficient to produce
noise-floor differences in the aggregated parameter trajectory
that compound over training.

We therefore interpret the runtime-magnitude difference as
\emph{implementation-level variance} that affects the
estimated effect size but not the direction. Both runtimes
produce a positive within-pair $\Delta$ with FedProx winning
the majority of paired seeds; the magnitude of the effect is
sensitive to the runtime choice in a way that is consistent
with high-variance noise rather than a systematic per-seed
mechanism.

\subsection{Summary}
\label{sec:engineered-summary}

On the engineered non-IID partition designed to create inter-
client drift on every minority class, FedProx outperforms
FedAvg by $+0.027$ macro-F1 (FedProx wins $9/10$ paired seeds)
under the pure-PyTorch reference implementation and by
$+0.007$ macro-F1 (FedProx wins $7/10$ paired seeds) under the
Flower simulation framework. The direction of effect is
consistent across both runtimes and is the opposite of the
IID-baseline null reported in
Section~\ref{sec:iid-baseline}, consistent with the
proximal-anchor mechanism activating when inter-client class
skew is present. The magnitude differs between runtimes by a
factor of approximately four, attributable to implementation-
level variance (floating-point non-associativity in parameter
aggregation, RNG routing through Ray scheduling, serialisation
across process boundaries) rather than to systematic per-seed
algorithmic differences. The pure-PyTorch result is treated as
the primary inferential result throughout the remainder of the
dissertation; the Flower result is treated as a cross-runtime
robustness check.

% ==============================================================================
\section{Per-Class Breakdown on the Engineered Partition}
\label{sec:engineered-per-class}
% ==============================================================================

The headline result in Section~\ref{sec:engineered-headline}
reports an aggregate $+0.027$ macro-F1 advantage for FedProx
over FedAvg on the engineered partition, averaged across the
seven DermaMNIST classes. The aggregate alone does not reveal
\emph{which} classes drive the improvement, and the proximal-
anchor mechanism makes a class-specific prediction that can be
tested directly: the regulariser should be active on classes
held by only a subset of clients (where inter-client drift
exists) and inert on classes held by every client (where no
drift exists). This section decomposes the headline by class
to test the prediction and to localise the clinical content of
the FedProx benefit.

\subsection{Setup}
\label{sec:engineered-per-class-setup}

The data are the same pure-PyTorch primary runs analysed in
Section~\ref{sec:engineered-headline}: ten paired seeds on the
engineered \texttt{balanced\_paired\_7\_clients} partition with
$E = 20$ local epochs, $R = 150$ communication rounds, $\mu =
0.01$ for FedProx and $\mu = 0$ for FedAvg. The per-class F1
scores reported below are computed on the same global test set
($n = 2{,}005$) at the per-seed best-validation checkpoint, and
are extracted directly from the \texttt{per\_class\_f1} field of
each result JSON. The paired-seed structure is preserved: within
each seed, FedProx and FedAvg share identical partition, initial
weights, and dataloader RNG state, so within-pair differences on
each class reflect only the proximal-term gating.

\subsection{Results}
\label{sec:engineered-per-class-results}

Table~\ref{tab:engineered-per-class} reports per-class test
F1 scores for FedAvg and FedProx across the ten paired seeds,
along with the within-pair $\Delta$ and the number of seeds on
which FedProx outperforms FedAvg.
Figure~\ref{fig:engineered-per-class} shows the per-class
validation F1 trajectories over 150 communication rounds.

\begin{table}[!htbp]
\centering
\caption{Per-class test F1 on the engineered partition
(pure-PyTorch primary, $n = 10$ paired seeds, mean $\pm$
standard deviation across seeds). The winning algorithm per
row is shown in bold; the $\Delta$ column is the mean within-
pair difference (FedProx $-$ FedAvg). Wins are the count of
paired seeds in which FedProx outperforms FedAvg on that class.}
\label{tab:engineered-per-class}
\small
\begin{tabular}{lcccc}
\toprule
Class (prevalence) & FedAvg & FedProx & $\Delta$ & Wins \\
\midrule
Actinic keratoses (3.27\%)        & $0.404 \pm 0.111$ & $\mathbf{0.471 \pm 0.040}$ & $+0.067$ & $9/10$ \\
Basal cell carcinoma (5.13\%)     & $\mathbf{0.475 \pm 0.046}$ & $0.473 \pm 0.027$ & $-0.001$ & $4/10$ \\
Benign keratosis (10.97\%)        & $0.440 \pm 0.033$ & $\mathbf{0.444 \pm 0.041}$ & $+0.004$ & $7/10$ \\
Dermatofibroma (1.15\%)           & $0.283 \pm 0.141$ & $\mathbf{0.296 \pm 0.070}$ & $+0.013$ & $4/10$ \\
Melanoma (11.11\%)                & $0.187 \pm 0.101$ & $\mathbf{0.301 \pm 0.026}$ & $\mathbf{+0.114}$ & $\mathbf{9/10}$ \\
Melanocytic nevi (67.05\%)        & $0.884 \pm 0.003$ & $\mathbf{0.886 \pm 0.004}$ & $+0.002$ & $7/10$ \\
Vascular lesions (1.41\%)         & $\mathbf{0.698 \pm 0.078}$ & $0.686 \pm 0.043$ & $-0.012$ & $5/10$ \\
\bottomrule
\end{tabular}
\end{table}

The per-class breakdown reveals a strongly non-uniform
distribution of the FedProx benefit across the seven classes.
The melanoma class stands out as the dominant contributor:
FedProx improves melanoma F1 by $+0.114$ on average, wins on
nine of ten paired seeds, and reduces the across-seed standard
deviation from $\pm 0.101$ (FedAvg) to $\pm 0.026$ (FedProx).
Actinic keratoses shows the next-largest improvement
($\Delta = +0.067$, $9/10$ wins). The remaining five classes
show small differences in either direction
($-0.012 \leq \Delta \leq +0.013$), all within their across-seed
standard deviations and not consistently signed across seeds.

The majority class \emph{melanocytic nevi} provides the
mechanism-cleanest observation: $\Delta = +0.002$ with both
algorithms reaching $F1 \approx 0.885$ and very small
across-seed standard deviations ($\pm 0.003$ for FedAvg,
$\pm 0.004$ for FedProx). Under the engineered partition,
melanocytic nevi is held in equal proportion by every one of
the seven clients, so no inter-client class skew exists on this
class for the proximal anchor to act on. The near-zero $\Delta$
and the tight across-seed agreement on the universal class is
the empirical signature predicted by the drift-control
mechanism: where there is no drift, FedProx produces the same
result as FedAvg, on every seed, on the same class. The
mel-nevi result is the within-experiment per-class control
introduced by the partition design.

The pattern visible in
Figure~\ref{fig:engineered-per-class} confirms the tabular
finding visually. The melanoma panel (bottom-left) shows the
largest persistent gap between the two algorithms across all
$150$ rounds: FedAvg plateaus at $\mathrm{F1} \approx 0.20$
while FedProx reaches $\mathrm{F1} \approx 0.30$, with the gap
visible from approximately round $50$ onwards. The melanocytic
nevi panel (bottom-centre) shows two curves that are essentially
indistinguishable across the entire training horizon. The other
five class panels show modest or noisy differences without a
visually persistent gap.

\begin{figure}[!htbp]
\centering
\includegraphics[width=\textwidth]{F_engineered_per_class.pdf}
\caption{Validation F1 over 150 communication rounds for each of
the seven DermaMNIST classes on the engineered partition (pure-
PyTorch primary, $n = 10$ paired seeds, mean $\pm$ SEM band).
Each panel shows the per-class trajectory for one class with the
class prevalence in parentheses and the within-pair test-set
$\Delta$ (FedProx $-$ FedAvg) annotated. The melanoma panel
(bottom-left) shows the largest persistent FedProx-FedAvg gap;
the melanocytic nevi panel (bottom-centre) shows two
indistinguishable trajectories, consistent with the
proximal-anchor mechanism being inert on a class held by every
client.}
\label{fig:engineered-per-class}
\end{figure}

\subsection{Mechanism interpretation}
\label{sec:engineered-per-class-mechanism}

The per-class pattern aligns with the engineered partition's
structural design and with the proximal-anchor mechanism's
prediction. The melanoma class is held by exactly two of the
seven clients (clients C4 and C5, the melanoma/vascular
specialty pair) and represents $11.1\%$ of the global training
distribution; under FedAvg, the two melanoma-holding clients
drift toward an over-specialised local optimum that
under-represents the other classes, and the aggregated update
incorporates this drift into the global model proportionally to
the clients' sample counts. The proximal anchor constrains how
far the two melanoma specialists can move per round, which
preserves a more balanced representation of melanoma in the
global model while still allowing the specialist clients to
exploit their local data. The melanocytic nevi class is held by
all seven clients in equal proportion ($670$--$673$ samples per
client), so all clients see the same melanocytic nevi
distribution locally and no inter-client drift exists on this
class for the proximal anchor to constrain --- yielding the
near-zero $\Delta$ observed empirically.

The contrast between the two clinically critical minority
classes \emph{melanoma} and \emph{actinic keratoses} (large
positive $\Delta$, mechanism-active) and the universal class
\emph{melanocytic nevi} (near-zero $\Delta$, mechanism-inert)
within the same experiment, on the same model, on the same
seeds, is the strongest single piece of per-class evidence for
the drift-control interpretation in this dissertation.

\subsection{Summary}
\label{sec:engineered-per-class-summary}

The $+0.027$ macro-F1 advantage reported in
Section~\ref{sec:engineered-headline} is class-localised: the
melanoma class contributes $+0.114$ macro-F1 in mean within-pair
$\Delta$ ($9/10$ paired wins), the actinic keratoses class
contributes $+0.067$ ($9/10$ wins), and the remaining five
classes contribute small differences within across-seed noise.
The mechanism-null prediction for the universal class
\emph{melanocytic nevi} is empirically confirmed
($\Delta = +0.002$, both algorithms within $\pm 0.005$ across
seeds). The class-by-class pattern is consistent with the
proximal anchor activating only where inter-client class skew
exists, as predicted by the engineered partition's structural
design.

% ==============================================================================
\section{Robustness Across Partition Designs}
\label{sec:partition-robustness}
% ==============================================================================

The headline result of Section~\ref{sec:engineered-headline} is
established on a single engineered partition that was designed
to maximise the proximal-anchor mechanism's effect. To check
that the FedProx-vs-FedAvg comparison reproduces beyond this
specific design --- and to contrast against the IID null
established in Section~\ref{sec:iid-baseline} --- we ran the
same paired-seed protocol on three alternative non-IID
partition designs covering different non-IID structures from
the federated-learning literature: a literature-standard
random-skew design (Dirichlet $\alpha = 0.1$), an extreme
singleton-specialist counterfactual, and the engineered
balanced-paired partition itself replicated on the Flower
runtime. All comparisons in this section share the Flower
runtime and the same ten paired seeds, so the only variable
between rows is the partition design.

\subsection{Partition designs}
\label{sec:partition-designs}

The IID and engineered (balanced-paired) partitions used as
reference rows in this section were introduced and visualised
earlier in the chapter
(Figures~\ref{fig:iid-partition}~and~\ref{fig:engineered-partition}).
Figure~\ref{fig:partition-comparison-dir-spec} introduces the
two new partition designs used here --- Dirichlet $\alpha = 0.1$
and the specialist (1-of-7) partition --- using the same
stacked-bar format and class palette so the structural
differences across all four partitions can be read by glancing
between the relevant figures. Per-panel $y$-axes are independent
because Dirichlet produces per-client totals that differ from
the other designs by more than an order of magnitude; the count
above each bar makes quantity skew explicit. Each design is
described in turn below.

\begin{figure}[!htbp]
\centering
\includegraphics[width=0.98\textwidth]{F_partition_comparison_dir_spec.pdf}
\caption{The two new partition designs introduced for the
robustness analysis (seed $42$). (a) Dirichlet $\alpha = 0.1$:
class-by-class proportions drawn from a Dirichlet, producing
severe joint label-and-quantity skew (per-client totals range
$91$--$3{,}721$). (b) Specialist (1-of-7): every minority class
is held by \emph{exactly one} client, with per-client sample
totals matched to the engineered partition
(Figure~\ref{fig:engineered-partition}) so that the only
variable between the engineered and specialist designs is the
structural change from paired to singleton minority-class
ownership.}
\label{fig:partition-comparison-dir-spec}
\end{figure}

\paragraph{Dirichlet $\alpha = 0.1$.}
The Dirichlet partitioning protocol of Hsu et al.\ (2019) is
the de-facto standard for synthetic non-IID benchmarking in
federated learning. For each class $c$ a $7$-client probability
vector $\boldsymbol{q}_c \sim \mathrm{Dir}_7(\alpha\boldsymbol{1})$ is
sampled, and the class-$c$ training samples are distributed to
the seven clients in proportions $q_{c,1}, \dots, q_{c,7}$. At
$\alpha = 0.1$ the resulting per-client distributions are
severely skewed: in the realisation used, client sample counts
range from $91$ to $3{,}721$ and individual classes are
concentrated on small subsets of clients. This combines label
skew with substantial quantity skew, which is the regime in
which most published non-IID benchmarks report FedAvg
degradation.

\paragraph{Specialist (1-of-7).}
The specialist partition assigns each minority class to
\emph{exactly one} client (rather than two, as in the
engineered design) while matching the per-client total sample
counts of the engineered partition exactly. Under this design
no inter-client drift exists on any minority class because each
minority class is held by only one client. The specialist
partition therefore tests the contrapositive of the engineered
design's mechanism prediction: if the proximal anchor's benefit
requires multiple clients to share a minority class, the
specialist partition should show a smaller per-class effect on
those minority classes than the engineered partition does.

\paragraph{Engineered balanced-paired (Flower).}
The same engineered partition used in
Section~\ref{sec:engineered-headline}, here run on the Flower
runtime so that all four partitions in this section share a
single runtime. This is the cross-runtime replication of the
primary result, included in this table for direct
within-runtime comparison.

\subsection{Results}
\label{sec:partition-robustness-results}

Table~\ref{tab:partition-robustness} reports test macro-F1
across the ten paired seeds on each partition. The IID baseline
from Section~\ref{sec:iid-baseline} is included as the reference
row. Figure~\ref{fig:partition-robustness} shows the per-round
validation macro-F1 trajectories on each partition.

\begin{table}[!htbp]
\centering
\caption{Robustness across four partition designs on the
Flower runtime, $n = 10$ paired seeds, all hyperparameters
held identical (only the partition function varies between
rows). Test macro-F1 reported as mean $\pm$ standard deviation
across seeds; the winning algorithm per row is shown in bold.
The IID row is the mechanism-null reference from
Section~\ref{sec:iid-baseline}.}
\label{tab:partition-robustness}
\small
\begin{tabular}{lccccc}
\toprule
Partition & FedAvg & FedProx & $\Delta$ & Wins \\
\midrule
IID                          & $\mathbf{0.5846 \pm 0.0197}$ & $0.5786 \pm 0.0190$ & $-0.006$ & $5/10$ \\
Dirichlet $\alpha = 0.1$     & $0.4796 \pm 0.0355$ & $\mathbf{0.4875 \pm 0.0341}$ & $+0.008$ & $6/10$ \\
Specialist (1-of-7)          & $0.4365 \pm 0.0236$ & $\mathbf{0.4446 \pm 0.0172}$ & $+0.008$ & $7/10$ \\
Engineered balanced-paired   & $0.4964 \pm 0.0150$ & $\mathbf{0.5033 \pm 0.0206}$ & $+0.007$ & $7/10$ \\
\bottomrule
\end{tabular}
\end{table}

The direction of the FedProx-vs-FedAvg effect is consistent
across the three non-IID partition designs (Dirichlet,
specialist, engineered): all three produce a positive mean
within-pair $\Delta$ in the range $+0.007$ to $+0.008$ macro-F1,
with FedProx winning the majority of paired seeds in each case
($6/10$, $7/10$, $7/10$). The IID partition produces the
opposite sign of small magnitude ($\Delta = -0.006$, FedProx
wins $5/10$). The four-way contrast between IID and the three
non-IID designs reproduces the qualitative pattern expected of
the proximal-anchor mechanism: positive direction under
non-IID, near-zero direction under IID. The magnitudes on
the three non-IID partitions are similar in size to the
across-seed standard deviations of either algorithm
($\sim 0.02$--$0.03$), so we report them as descriptive
observations of direction rather than as quantitatively
distinguishable effect sizes.

The absolute macro-F1 levels vary substantially across
partitions, reflecting differences in task difficulty: IID is
the easiest setting ($\sim 0.58$ macro-F1) because every client
sees the global class distribution proportionally; the
specialist partition is the hardest ($\sim 0.44$) because each
minority class is held by only one client and the global model
must aggregate over highly specialised local optima; the
Dirichlet $\alpha = 0.1$ and engineered partitions sit between
the two extremes. The within-pair $\Delta$ comparison in
Table~\ref{tab:partition-robustness} is paired-seed within each
row and therefore controls for these absolute-level differences
across partitions.

\begin{figure}[!htbp]
\centering
\includegraphics[width=\textwidth]{F_partition_robustness.pdf}
\caption{Validation macro-F1 over $150$ communication rounds
on four DermaMNIST partition designs (Flower runtime, $n = 10$
paired seeds, mean $\pm$ SEM band). The IID panel (top-left)
shows the two algorithms' trajectories essentially
superimposed; the three non-IID panels (top-right, bottom-left,
bottom-right) show a small but visible persistent gap with
FedProx above FedAvg from approximately round~$30$ onwards.}
\label{fig:partition-robustness}
\end{figure}

\subsection{Cross-design observation}
\label{sec:partition-cross-design}

The three non-IID partitions tested differ substantially in
their structural design: the Dirichlet $\alpha = 0.1$ partition
produces random and severe class skew with large per-client
quantity differences (client sample counts $91$ to $3{,}721$);
the specialist partition assigns each minority class to exactly
one client; the engineered partition assigns each minority
class to exactly two co-training clients with mild quantity
skew. Despite these substantial structural differences, the
three partitions produce within-pair $\Delta$ values within
$\pm 0.001$ macro-F1 of each other ($+0.007$, $+0.008$,
$+0.008$). At the chosen sample size of $n = 10$ paired seeds,
the across-seed noise dominates any potential structural
difference between non-IID designs at the aggregate macro-F1
level.

A more discriminating cross-design observation emerges at the
per-class level. On the melanoma class, the engineered
partition produces $\Delta_{\mathrm{melanoma}} = +0.068$ ($7/10$
paired wins) while the specialist partition produces
$\Delta_{\mathrm{melanoma}} = -0.015$ ($4/10$ wins). The
structural explanation aligns with the proximal-anchor
mechanism: under the engineered partition, melanoma is held by
two clients (the melanoma/vascular specialty pair), which
drift apart during local SGD and require reconciliation at
aggregation --- the regime in which the proximal anchor is
mechanically active. Under the specialist partition, melanoma
is held by only one client, so no inter-client drift exists on
the melanoma class and the proximal anchor is structurally
inert for that class. The per-class melanoma contrast between
engineered and specialist therefore corroborates the
mechanism-based interpretation of the per-class breakdown in
Section~\ref{sec:engineered-per-class}: the proximal anchor's
contribution is conditional on multiple clients sharing the
class.

\subsection{Summary}
\label{sec:partition-robustness-summary}

Across three structurally different non-IID partition designs
(Dirichlet $\alpha = 0.1$, specialist, engineered
balanced-paired), the within-pair $\Delta$ is positive and
similar in magnitude ($+0.007$ to $+0.008$ macro-F1, FedProx
wins $6$--$7$ of $10$ paired seeds in each case). Under the
IID partition the within-pair $\Delta$ is near-zero and
negative ($-0.006$). The four-way contrast reproduces the
direction expected of the proximal-anchor mechanism while
revealing that, at $n = 10$ paired seeds and at the aggregate
macro-F1 level, the choice of non-IID partition design does
not strongly modulate the magnitude of the FedProx benefit.
A more discriminating per-class contrast --- melanoma under
engineered vs specialist --- corroborates the mechanism-based
interpretation that the proximal anchor is active when
multiple clients share a minority class and inert when only
one client holds it.

% ==============================================================================
\section{Direct Mechanism Evidence: Per-Client Update Norms}
\label{sec:update-norms}
% ==============================================================================

The results reported so far in this chapter argue for the
proximal-anchor mechanism \emph{indirectly}, by interpreting
patterns in test macro-F1 (cross-partition direction
consistency, per-class melanoma vs mel-nevi contrast, IID
null). A more direct line of evidence is available: because the
proximal regulariser $(\mu/2)\lVert w - w^t \rVert^2$ is a
penalty on the parameter increment, FedProx is expected to
produce systematically smaller per-client parameter updates
than FedAvg when client drift is present. This section reports
this measurement directly, using per-round per-client update-
norm logging from the Flower runtime sweeps.

\subsection{What is measured}
\label{sec:update-norms-setup}

For each communication round $t$ and each client $i$, we log
the Euclidean distance between that client's locally-updated
parameter vector at the end of the round and the global
parameter vector broadcast at the start of the round:
\[
\rho_{t,i} \;=\; \lVert w_i^{t+1} - w^t \rVert_2.
\]
This quantity has a clean mechanism interpretation: under
FedAvg, $\rho_{t,i}$ is determined entirely by the local
training trajectory; under FedProx, the proximal regulariser
explicitly penalises $\lVert w - w^t \rVert$ during local
training, so $\rho_{t,i}$ is expected to be smaller. The
proximal anchor's drift-control effect, if present, should
appear directly in this measurement, independent of any
outcome-level metric.

The Flower runtime client-side training code logs $\rho_{t,i}$
for every (round, client) pair when the
\texttt{--log-update-norms} flag is enabled. This was the case
for the Flower engineered-partition baseline sweep
(\texttt{flower\_C0\_baseline}) and the Flower IID baseline
sweep (\texttt{flower\_C0\_iid\_baseline}); the remaining
Flower statistical-heterogeneity sweeps used in
Section~\ref{sec:partition-robustness} do not have update-norm
logging in their result directories. We therefore report
update-norm results on the engineered and IID partitions, the
two endpoints of the heterogeneity-intensity axis.

\subsection{Results}
\label{sec:update-norms-results}

Table~\ref{tab:update-norms} reports the mean update norm
across all (round, client) pairs, averaged across paired seeds,
for FedAvg and FedProx on each partition.
Figure~\ref{fig:update-norms} shows the per-round trajectory of
the mean update norm across the 150 communication rounds.

\begin{table}[!htbp]
\centering
\caption{Per-client update norm
$\lVert w_i^{t+1} - w^t \rVert_2$ averaged across all (round,
client) pairs and across paired seeds (Flower runtime). The
ratio column reports the mean FedAvg norm divided by the mean
FedProx norm for each partition.}
\label{tab:update-norms}
\small
\begin{tabular}{lccc}
\toprule
Partition & FedAvg $\lVert \Delta w \rVert$ & FedProx $\lVert \Delta w \rVert$ & Ratio (FA / FP) \\
\midrule
Engineered (non-IID)   & $1.523$ & $\mathbf{1.029}$ & $\mathbf{1.48\times}$ \\
IID (mechanism null)   & $0.825$ & $\mathbf{0.724}$ & $1.14\times$ \\
\bottomrule
\end{tabular}
\end{table}

Two observations follow from Table~\ref{tab:update-norms}.

First, on both partitions FedProx produces smaller per-client
update norms than FedAvg --- the direction predicted by the
proximal regulariser. This is the direct empirical signature
of the drift-control mechanism: when the proximal regulariser
is active, per-round client updates are bounded in magnitude.

Second, the magnitude of the FedAvg/FedProx ratio differs
substantially between the two partitions. On the engineered
non-IID partition the FedAvg update norm is $1.48\times$ that of
FedProx ($1.523$ vs $1.029$); on the IID partition the ratio
shrinks to $1.14\times$ ($0.825$ vs $0.724$). The same
proximal coefficient $\mu = 0.01$ produces a substantially
stronger empirical drift-control effect on the non-IID
partition than on the IID partition. This is mechanism-
consistent: under IID, all clients are drifting toward
similar local optima (because their local objectives are
i.i.d.\ samples of the same global distribution), so there is
little client drift to constrain and the proximal anchor's
empirical effect is modest. Under the engineered non-IID
partition, clients drift toward distinct specialised local
optima, and the proximal anchor's effect is substantially
larger.

\begin{figure}[!htbp]
\centering
\includegraphics[width=\textwidth]{F_update_norms.pdf}
\caption{Per-client update norm
$\lVert w_i^{t+1} - w^t \rVert_2$ across $150$ communication
rounds on the Flower runtime, averaged across clients within
each round and across paired seeds (mean $\pm$ SEM band). On
the engineered non-IID partition (left), FedProx maintains a
visibly smaller per-client update norm than FedAvg from
round~$1$ onwards; the gap persists across the entire $150$-
round horizon. On the IID partition (right), both algorithms'
update norms decay rapidly toward zero by round~$\sim 80$ and
the FedProx-FedAvg gap is correspondingly smaller. The cross-
partition contrast shows the proximal anchor's empirical effect
scales with the amount of inter-client drift present in the
training setup.}
\label{fig:update-norms}
\end{figure}

The convergence-time profile of the update norms also differs
between partitions. On the IID partition, both algorithms'
update norms decay to near zero by round $\sim 80$, indicating
that all seven clients have approximately converged to the
same global optimum and per-round changes are minimal. On the
engineered partition, the update norms decay more slowly and
remain non-trivial across the entire $150$-round horizon ---
clients continue to make meaningful per-round parameter
changes even in late training because their local optima
differ from the global aggregate. The FedAvg/FedProx ratio
remains stable across the late-training horizon on the
engineered partition rather than collapsing toward one,
indicating that the proximal regulariser remains active
throughout training rather than becoming inert at convergence.

\subsection{Trajectory-volatility corroboration}
\label{sec:update-norms-volatility}

An outcome-level signal that corroborates the parameter-level
finding is the per-round volatility of the validation
macro-F1 trajectory. Smaller per-client parameter updates
should produce smoother round-to-round changes in the global
model and therefore smoother validation curves; larger
updates should produce more volatile trajectories. We
quantify this by the mean absolute round-to-round change
in validation macro-F1 across the $150$ communication rounds,
averaged across the ten paired seeds of the pure-PyTorch
primary sweep:
\[
\text{volatility}_a \;=\;
\mathrm{mean}_r \big|\mathrm{F1}_a(r) - \mathrm{F1}_a(r-1)\big|.
\]
On the pure-PyTorch engineered headline,
$\text{volatility}_{\mathrm{FedAvg}} = 0.0414$ and
$\text{volatility}_{\mathrm{FedProx}} = 0.0225$, with the
FedAvg/FedProx ratio of $1.84\times$. Crucially, the ratio
exceeds one on every one of the ten paired seeds (per-seed
range: $1.30\times$ to $2.95\times$), so this is not a feature
of any single seed. FedProx's smoother validation trajectory
is a downstream consequence of its smaller per-client update
norm.

\subsection{Summary}
\label{sec:update-norms-summary}

Direct measurement of per-client parameter updates confirms
the proximal-anchor mechanism empirically. On the engineered
non-IID partition, the mean FedAvg per-client update norm is
$1.48\times$ that of FedProx; on the IID partition, the ratio
shrinks to $1.14\times$. The same proximal coefficient
produces a substantially larger empirical effect on parameter
updates when more inter-client drift is present, matching the
mechanism prediction that the regulariser acts on client
drift. The outcome-level signature of this parameter-level
effect is a $1.84\times$ smoother validation-macro-F1
trajectory under FedProx than under FedAvg on the engineered
partition, holding on every paired seed.

% ==============================================================================
\section{Federation Tax}
\label{sec:federation-tax}
% ==============================================================================

To contextualise the within-federation FedAvg-vs-FedProx
comparison against the absolute scale set by
Section~\ref{sec:centralised-baseline}, we express each
federated configuration's mean test macro-F1 as a percentage of
the centralised reference ($0.5600$ macro-F1). The four
federated configurations all use the engineered
\texttt{balanced\_paired\_7\_clients} partition, $n = 10$ paired
seeds, and identical hyperparameters.

\begin{table}[!htbp]
\centering
\caption{Federation tax. Each row's recovery percentage is its
test macro-F1 mean divided by the centralised baseline mean of
$0.5600$.}
\label{tab:federation-tax}
\small
\begin{tabular}{lcc}
\toprule
Configuration & Test macro-F1 & Recovery vs centralised \\
\midrule
Centralised (non-federated)         & $\mathbf{0.5600 \pm 0.0240}$ & $100.0\%$ \\
FedAvg, pure-PyTorch                & $0.4814 \pm 0.0254$ & $86.0\%$ \\
FedProx, pure-PyTorch               & $0.5081 \pm 0.0140$ & $90.7\%$ \\
FedAvg, Flower                      & $0.4964 \pm 0.0150$ & $88.6\%$ \\
FedProx, Flower                     & $0.5033 \pm 0.0206$ & $89.9\%$ \\
\bottomrule
\end{tabular}
\end{table}

The federation tax --- the gap between centralised training
and any federated configuration --- is $9$--$14$ percentage
points of macro-F1, substantially larger than the within-
federation FedProx-vs-FedAvg improvement ($4.7$ percentage
points on pure-PyTorch, $1.3$ on Flower). The thesis result
is therefore best framed not as ``federated training matches
centralised'' but as ``federated FedProx closes a modest
portion of the federation gap that federated FedAvg leaves
open.'' The remaining $9$--$10$ percentage-point gap to
centralised is the cost of federation itself and is not
addressed by the choice of FL algorithm.

% ==============================================================================
\section{System Heterogeneity: Motivation and Scope}
\label{sec:system-het-motivation}
% ==============================================================================

The chapter up to this point has focused on \emph{statistical}
heterogeneity --- differences in the data distribution held by
each client. Federated-learning theory and practice identifies
a second canonical axis of heterogeneity, equally important in
production deployments: \emph{system} heterogeneity, in which
clients differ in the amount of local computation they can
contribute per communication round. In real federated
deployments, institutions run on heterogeneous hardware, have
different on-call training windows, and cannot all complete the
same amount of local work within a round's wall-clock deadline.
Some clients (``stragglers'') return partial updates after fewer
local optimisation steps; the question is how the FL algorithm
should handle these partial updates.

The proximal term in FedProx was originally motivated for this
regime. Li et al.\ (2020), Theorem~4, proves a convergence
guarantee for FedProx under \emph{$\gamma$-inexact} local
updates: each client may submit a parameter vector
$w_i^{t+1}$ that does not fully minimise its local objective,
provided the residual inexactness is bounded. The proximal
regulariser $(\mu/2)\lVert w - w^t \rVert^2$ constrains how far
an inexact local update can drift from the broadcast global
anchor $w^t$, making partial work from stragglers usable for
aggregation without destabilising the global model. FedAvg has
no such guarantee, and its size-weighted aggregation rule
treats partial-work and full-work updates symmetrically,
weighted only by dataset size.

This chapter tests whether the FedProx-vs-FedAvg pattern
established under statistical heterogeneity in the preceding
sections persists, attenuates, or amplifies when clients have
heterogeneous per-round local-epoch budgets. The same
engineered \texttt{balanced\_paired\_7\_clients} partition,
the same ten paired seeds, and the same Flower simulation
runtime are used throughout, so that the only variable
between this chapter's experiments and the engineered
non-IID baseline of Section~\ref{sec:engineered-headline} is
the per-(round, client) local-epoch allocation. An IID
cross-partition control mirrors the contrast structure used in
Section~\ref{sec:partition-robustness}.

\paragraph{Scope and terminology.}
We use ``system heterogeneity'' as shorthand for what is
strictly \emph{computational-work heterogeneity}: variability
in the number of local SGD steps each client performs per
round. We do not simulate hardware speed, network latency, or
device availability; the manipulation acts only on the
algorithm via heterogeneous per-(round, client) local-epoch
budgets $\tau_i$. We restrict scope to the FedAvg/FedProx
comparison; alternative algorithms designed specifically for
heterogeneous local-step regimes are out of scope for this
chapter.

% ==============================================================================
\section{Experimental Design and Conditions}
\label{sec:system-het-design}
% ==============================================================================

Three system-heterogeneity conditions are evaluated, denoted
$C_0$, $C_1$, and $C_2$. Each condition defines a per-(round,
client) local-epoch schedule that is fully determined by the
random seed; paired FedAvg/FedProx runs at the same seed
therefore see an identical schedule and the within-pair
$\Delta$ reflects only the algorithmic gating of the proximal
term. The schedules under seed~$42$ are visualised in
Figure~\ref{fig:sh-schedules}.

\paragraph{C0 --- no system heterogeneity (baseline).}
Every client performs $E = 20$ local epochs every round, for
all $150$ communication rounds. This is the same setting used
in the engineered-partition headline of
Section~\ref{sec:engineered-headline} and provides the
within-runtime within-partition baseline against which the
two system-heterogeneity conditions are compared. In the
schedule heatmap (Figure~\ref{fig:sh-schedules}, left), every
(round, client) cell is yellow, indicating the maximum
local-epoch budget.

\paragraph{C1 --- fixed stragglers.}
Two designated clients (C5 and C6) always perform
$E_{\text{straggler}} = 5$ local epochs; the remaining five
clients always perform the full $E = 20$. This represents the
deployment pattern in which specific institutions consistently
have weaker compute infrastructure or shorter training windows.
The schedule heatmap (Figure~\ref{fig:sh-schedules}, centre)
shows a solid vertical bar of $E = 5$ on columns C5 and C6
across all $150$ rounds, with the remaining columns at
$E = 20$.

\paragraph{Caveat on the C1 design.} On the engineered
\texttt{balanced\_paired\_7\_clients} partition, client~C5 is
part of the melanoma/vascular specialty pair (the pair that
drives the per-class melanoma advantage in
Section~\ref{sec:engineered-per-class}), and client~C6 is the
melanocytic-nevi-only generalist. C1 therefore couples the
``system heterogeneity'' manipulation with a specific clinical
role assignment under the engineered partition. The C2
condition below does not have this confound (the straggler
identity rotates randomly per round), so C2 is the cleaner
test of the pure system-heterogeneity effect. C1 is retained
for descriptive realism and as a contrast to C2.

\paragraph{C2 --- random stragglers.}
Each communication round, $\lfloor 0.5 \times 7 \rfloor + 1 = 4$
clients are randomly designated stragglers (sampled without
replacement). Each straggler performs $E_i \sim
\mathrm{Uniform}\{1, 2, \dots, 19\}$ local epochs; the
remaining three clients perform the full $E = 20$. Aggregated
across the ten paired seeds, every client serves as a
straggler in approximately $57\%$ of rounds. The schedule
heatmap (Figure~\ref{fig:sh-schedules}, right) shows the
characteristic salt-and-pepper pattern of randomised straggler
assignment, with mean per-cell value of approximately $14.1$
epochs (mixture of full-work cells at $E = 20$ and partial-work
cells at $E_i \in \{1, \dots, 19\}$). C2 follows the protocol
of Li et al.\ (2020), §5.2 up to two implementation details:
stragglers always perform at least one local optimisation step
(the Li et al.\ protocol allows $E_i = 0$), and both algorithms
include partial updates in aggregation. We replicate the
canonical Li et al.\ asymmetric setup (FedAvg drops stragglers
while FedProx includes them) separately in
Section~\ref{sec:sh-asymmetric}.

\begin{figure}[!htbp]
\centering
\includegraphics[width=\textwidth]{F_sh_schedules.pdf}
\caption{Per-(round, client) local-epoch schedules under the
three system-heterogeneity conditions, reconstructed
deterministically from seed~$42$. Each heatmap row is a
communication round (rounds $1$--$150$, bottom to top) and
each column is one of the seven clients; cell colour shows the
local-epoch budget $E$ assigned to that (round, client) pair
on a $0$--$20$ scale. \textbf{Left (C0):} every client performs
the maximum $E = 20$ in every round (uniform yellow).
\textbf{Centre (C1):} clients C5 and C6 always perform
$E = 5$ (solid dark band on the two rightmost columns); the
remaining five clients always perform $E = 20$.
\textbf{Right (C2):} four of seven clients per round are
randomly designated stragglers with $E_i$ uniformly drawn from
$\{1, \dots, 19\}$; the remaining three clients perform the
full $E = 20$.}
\label{fig:sh-schedules}
\end{figure}

\paragraph{Training setup (held constant across conditions).}
All hyperparameters other than the schedule are held identical
to the engineered-partition headline of
Section~\ref{sec:engineered-headline}: same
\texttt{DermMNISTCNN} architecture, SGD with learning rate
$0.01$, momentum $0.9$, batch size $32$, weight decay $0$,
$R = 150$ communication rounds, full participation $C = 1.0$,
FedProx proximal coefficient $\mu = 0.01$, and the same ten
paired seeds. The Flower simulation runtime is used throughout
this chapter so that all C0/C1/C2 results share a single
runtime and the within-condition $\Delta$ does not mix
algorithmic effect with cross-runtime variation.

% ==============================================================================
\section{Results on the Engineered Partition}
\label{sec:sh-engineered-results}
% ==============================================================================

This section reports the FedAvg-vs-FedProx paired comparison
under each of the three system-heterogeneity conditions on the
engineered partition. The same ten paired seeds are used across
all conditions; within each seed, FedAvg and FedProx share an
identical local-epoch schedule, partition assignment, and
initial weights, so within-pair differences reflect only the
algorithmic gating of the proximal term.

\subsection{Results}
\label{sec:sh-engineered-table}

Table~\ref{tab:sh-engineered} reports test macro-F1 at the
best-validation checkpoint across the ten paired seeds under
each condition. Figure~\ref{fig:sh-engineered-convergence}
shows the corresponding per-round validation macro-F1
trajectories.

\begin{table}[!htbp]
\centering
\caption{System-heterogeneity results on the engineered
partition (Flower runtime, $n = 10$ paired seeds). Test
macro-F1 reported as mean $\pm$ standard deviation across
seeds; the winning algorithm per row is shown in bold. The
$\Delta$ column is the mean within-pair difference (FedProx
$-$ FedAvg). The C0 baseline row reproduces the engineered-
partition Flower result from
Section~\ref{sec:engineered-headline} for direct comparison.}
\label{tab:sh-engineered}
\small
\begin{tabular}{lccccc}
\toprule
Condition & FedAvg & FedProx & $\Delta$ & Wins \\
\midrule
C0 (no system het)                & $0.4964 \pm 0.0150$ & $\mathbf{0.5033 \pm 0.0206}$ & $+0.007$ & $7/10$ \\
C1 (fixed stragglers, C5+C6 @ E=5) & $0.4978 \pm 0.0163$ & $\mathbf{0.5111 \pm 0.0196}$ & $+0.013$ & $8/10$ \\
C2 (random stragglers, 4/7)        & $0.4928 \pm 0.0277$ & $\mathbf{0.5102 \pm 0.0168}$ & $+0.017$ & $7/10$ \\
\bottomrule
\end{tabular}
\end{table}

\begin{figure}[!htbp]
\centering
\includegraphics[width=\textwidth]{F_sh_engineered_convergence.pdf}
\caption{Validation macro-F1 over $150$ communication rounds
on the engineered \texttt{balanced\_paired\_7\_clients}
partition under each of the three system-heterogeneity
conditions (Flower runtime, $n = 10$ paired seeds, mean $\pm$
SEM band). The FedProx curve (dark slate blue) maintains a
visible upward separation from the FedAvg curve (mint green)
in all three conditions, with the gap widening from C0 (left)
to C2 (right).}
\label{fig:sh-engineered-convergence}
\end{figure}

\subsection{Interpretation}
\label{sec:sh-engineered-interpretation}

The within-pair $\Delta$ rises monotonically with straggler
severity across the three conditions: $+0.007$ macro-F1 at the
baseline C0, $+0.013$ at C1 with two consistently slow clients,
and $+0.017$ at C2 with four randomly designated stragglers
per round. FedProx wins the majority of paired seeds in each
condition ($7/10$, $8/10$, $7/10$). The convergence figure
shows this pattern visually: the FedAvg-vs-FedProx gap is
narrowest under C0 and visibly wider under C1 and C2,
particularly in the late-training plateau.

The monotonic increase aligns with the canonical motivation
of Li et al.\ (2020) for the proximal regulariser. As more
local computation per client becomes inexact ($E_i < E_{\max}$
on a growing fraction of (round, client) cells across the
three conditions), the parameter-update trajectories of the
participating clients diverge more from the broadcast global
anchor $w^t$. The proximal anchor's drift-control role is
therefore expected to contribute more under increasing
straggler severity, and the empirical monotonic pattern in
Table~\ref{tab:sh-engineered} is consistent with this
mechanism.

The absolute magnitudes are small in macro-F1 units
($0.007$--$0.017$), comparable in size to the across-seed
standard deviation of either algorithm under each condition
($\pm 0.015$--$\pm 0.028$). We therefore report the
condition-level $\Delta$ values descriptively, as a
qualitatively consistent direction-of-effect pattern rather
than as quantitatively distinguishable effect sizes between
adjacent conditions. The direction-of-effect pattern (positive
under non-IID stragglers, monotonic with severity) is the
substantive observation.

The C1 condition's positive $\Delta$ should be read together
with the partition-design caveat noted in
Section~\ref{sec:system-het-design}: under the engineered
partition, the two designated C1 stragglers (C5 and C6) are
the melanoma/vascular specialist client and the nevi-only
generalist client respectively. The C1 result therefore
reflects the specific scenario in which a clinical specialist
client and the generalist client are both consistently slow.
The C2 condition rotates straggler identity randomly per
round and therefore does not have this confound; the $+0.017$
macro-F1 advantage under C2 is the cleaner direct test of the
system-heterogeneity effect in the engineered-partition
configuration.

\subsection{Connection to the parameter-level evidence}
\label{sec:sh-engineered-connection}

The outcome-level pattern reported in
Table~\ref{tab:sh-engineered} is consistent with the direct
parameter-level mechanism evidence established in
Section~\ref{sec:update-norms}. In particular, the per-client
update-norm ratio $\lVert \Delta w \rVert_{\mathrm{FedAvg}} /
\lVert \Delta w \rVert_{\mathrm{FedProx}}$ on the engineered
partition was $1.48\times$ at the C0 baseline. The
system-heterogeneity conditions, in which local-epoch counts
$E_i$ vary across (round, client) cells, are precisely the
regime in which the proximal anchor's drift-control role is
theoretically most useful (Li et al.\ 2020, Theorem~4); the
larger within-pair $\Delta$ values observed under C1 and C2
relative to C0 are the outcome-level manifestation of this
broader drift-control role.

\subsection{Summary}
\label{sec:sh-engineered-summary}

Under the three system-heterogeneity conditions on the
engineered partition, the FedProx-vs-FedAvg within-pair
$\Delta$ rises monotonically with straggler severity (C0:
$+0.007$, C1: $+0.013$, C2: $+0.017$ macro-F1), and FedProx
wins the majority of paired seeds in each condition. The
absolute magnitudes are small but the direction-of-effect
pattern is consistent with the canonical proximal-anchor
mechanism: more inexact local work produces more
opportunity for the drift-control term to act, and the
empirical $\Delta$ increases correspondingly. The
randomised-straggler C2 condition is the cleanest test (no
clinical-role confounding) and produces the largest
within-pair advantage.

% ==============================================================================
\section{Per-Class Breakdown under C2}
\label{sec:sh-c2-per-class}
% ==============================================================================

The aggregate macro-F1 result reported in
Section~\ref{sec:sh-engineered-results} averages over the seven
DermaMNIST classes and does not reveal which classes contribute
to the $+0.017$ within-pair $\Delta$ under the C2 random-
straggler condition. This section decomposes the C2 result by
class, mirroring the per-class breakdown performed for the
no-system-het engineered baseline in
Section~\ref{sec:engineered-per-class}. The expected pattern,
based on the partition's structural design, is that classes
held by exactly two clients should show a positive $\Delta$
(inter-client drift exists for the proximal anchor to act on)
while the majority class melanocytic nevi (held by every
client) should show $\Delta \approx 0$.

\subsection{Setup}
\label{sec:sh-c2-per-class-setup}

The data are the same C2 paired runs analysed in
Section~\ref{sec:sh-engineered-results}: ten paired seeds on
the engineered \texttt{balanced\_paired\_7\_clients} partition
with the random-straggler schedule, $\mu = 0.01$ for FedProx
and $\mu = 0$ for FedAvg, all other hyperparameters identical
to the C0 baseline. Per-class F1 scores are computed on the
global test set ($n = 2{,}005$ samples) at each run's best-
validation checkpoint.

\subsection{Results}
\label{sec:sh-c2-per-class-results}

Table~\ref{tab:sh-c2-per-class} reports per-class test F1
across the ten paired seeds. Figure~\ref{fig:sh-c2-per-class}
shows the corresponding per-class validation F1 trajectories.

\begin{table}[!htbp]
\centering
\caption{Per-class test F1 under the C2 random-straggler
condition on the engineered partition (Flower runtime, $n = 10$
paired seeds, mean $\pm$ standard deviation across seeds).
The winning algorithm per row is shown in bold; the $\Delta$
column is the mean within-pair difference (FedProx $-$ FedAvg).
Wins are the count of paired seeds in which FedProx outperforms
FedAvg on that class.}
\label{tab:sh-c2-per-class}
\small
\begin{tabular}{lcccc}
\toprule
Class (prevalence) & FedAvg & FedProx & $\Delta$ & Wins \\
\midrule
Actinic keratoses (3.27\%)        & $0.454 \pm 0.060$ & $\mathbf{0.481 \pm 0.038}$ & $+0.026$ & $7/10$ \\
Basal cell carcinoma (5.13\%)     & $0.492 \pm 0.057$ & $\mathbf{0.497 \pm 0.023}$ & $+0.005$ & $4/10$ \\
Benign keratosis (10.97\%)        & $0.428 \pm 0.066$ & $\mathbf{0.432 \pm 0.045}$ & $+0.003$ & $4/10$ \\
Dermatofibroma (1.15\%)           & $0.255 \pm 0.065$ & $\mathbf{0.288 \pm 0.103}$ & $+0.034$ & $6/10$ \\
Melanoma (11.11\%)                & $0.233 \pm 0.065$ & $\mathbf{0.318 \pm 0.058}$ & $\mathbf{+0.085}$ & $\mathbf{10/10}$ \\
Melanocytic nevi (67.05\%)        & $0.884 \pm 0.005$ & $\mathbf{0.886 \pm 0.004}$ & $+0.001$ & $7/10$ \\
Vascular lesions (1.41\%)         & $\mathbf{0.703 \pm 0.066}$ & $0.670 \pm 0.039$ & $-0.033$ & $4/10$ \\
\bottomrule
\end{tabular}
\end{table}

The per-class decomposition under C2 closely mirrors the
per-class structure observed in the no-system-het baseline of
Section~\ref{sec:engineered-per-class}. Melanoma stands out as
the dominant contributor: FedProx improves melanoma F1 by
$+0.085$ on average and wins on \emph{every} one of the ten
paired seeds. The actinic keratoses class shows the next-
largest positive difference ($\Delta = +0.026$, $7/10$ wins),
followed by dermatofibroma ($\Delta = +0.034$, $6/10$ wins).
The majority class melanocytic nevi again produces a
near-zero $\Delta$ ($+0.001$ macro-F1, both algorithms within
$\pm 0.005$ across seeds) --- the expected signature of the
proximal-anchor mechanism on a class held by every client.

\begin{figure}[!htbp]
\centering
\includegraphics[width=\textwidth]{F_sh_c2_per_class.pdf}
\caption{Validation F1 over $150$ communication rounds for each
of the seven DermaMNIST classes under the C2 random-straggler
condition on the engineered partition (Flower runtime, $n = 10$
paired seeds, mean $\pm$ SEM band). Each panel shows the per-
class trajectory for one class with the class prevalence in
parentheses and the within-pair test-set $\Delta$ (FedProx $-$
FedAvg) annotated. The melanoma panel (bottom-left) shows the
largest persistent FedProx-vs-FedAvg gap; the melanocytic nevi
panel (bottom-centre) shows two indistinguishable trajectories.}
\label{fig:sh-c2-per-class}
\end{figure}

\subsection{Comparison with the C0 baseline}
\label{sec:sh-c2-per-class-vs-c0}

The C0 (no system het) baseline reported in
Section~\ref{sec:engineered-per-class} showed a melanoma
$\Delta$ of $+0.114$ on the pure-PyTorch primary; under the
Flower runtime which is used in this chapter the corresponding
C0 melanoma $\Delta$ is $+0.068$ ($7/10$ wins). The C2
random-straggler condition produces a melanoma $\Delta$ of
$+0.085$ ($10/10$ wins) --- larger in mean magnitude than the
Flower C0 baseline and unanimous in direction across seeds.
The melanocytic-nevi $\Delta$ remains near zero in both
conditions ($+0.002$ at C0 Flower, $+0.001$ at C2). The
per-class pattern that the macro-F1 advantage is concentrated
on the melanoma class and absent on the universal melanocytic
nevi class is preserved under system heterogeneity, with the
melanoma signal becoming slightly stronger and more
seed-consistent under C2 than under C0.

\subsection{Summary}
\label{sec:sh-c2-per-class-summary}

The $+0.017$ macro-F1 advantage reported under C2 in
Section~\ref{sec:sh-engineered-results} is class-localised
in the same way as the no-system-het engineered headline:
melanoma contributes a within-pair $\Delta$ of $+0.085$ with
FedProx winning all $10$ paired seeds, and the universal class
melanocytic nevi shows $\Delta \approx 0$ as predicted by the
drift-control mechanism. The per-class structure of the
FedProx advantage is therefore preserved under the random-
straggler schedule, with the melanoma class --- the clinically
most important diagnosis in this dataset --- as the dominant
contributor to the aggregate macro-F1 effect.

% ==============================================================================
\section{Replication of Li et al.\ (2020) §5.2 Asymmetric Protocol}
\label{sec:sh-asymmetric}
% ==============================================================================

The C0/C1/C2 results reported in
Sections~\ref{sec:sh-engineered-results}--\ref{sec:sh-c2-per-class}
use a \emph{symmetric} inclusion policy: both FedAvg and FedProx
aggregate over all clients including stragglers, weighted by
dataset size. The canonical straggler protocol of
Li et al.\ (2020), §5.2 is asymmetric: FedAvg \emph{drops}
straggler updates (clients that do not complete the full
$E_{\max}$ local epochs are excluded from aggregation), while
FedProx \emph{includes} their partial updates as
$\gamma$-inexact contributions. To verify that the canonical
asymmetric setup reproduces qualitatively on DermaMNIST, we
additionally run the FedAvg-drop arm on the engineered
partition under the C2 random-straggler schedule. Combined with
the existing FedProx-include arm from the symmetric C2 data of
Section~\ref{sec:sh-engineered-results} on the same ten seeds,
this gives the paired comparison in
Table~\ref{tab:sh-asymmetric}.

\begin{table}[!htbp]
\centering
\caption{Replication of the Li et al.\ (2020) §5.2 asymmetric
straggler protocol on the engineered partition with C2 random
stragglers (Flower runtime, $n = 10$ paired seeds). Test
macro-F1 reported as mean $\pm$ standard deviation across
seeds; the winning configuration is shown in bold.}
\label{tab:sh-asymmetric}
\small
\begin{tabular}{lccc}
\toprule
Configuration & Test macro-F1 & $\Delta$ & Wins \\
\midrule
FedAvg (drops straggler updates)         & $0.4423 \pm 0.0192$ & ---       & ---           \\
FedProx (includes straggler updates)     & $\mathbf{0.5102 \pm 0.0168}$ & $\mathbf{+0.068}$ & $\mathbf{10/10}$ \\
\bottomrule
\end{tabular}
\end{table}

The FedProx-include configuration outperforms FedAvg-drop by a
mean within-pair difference of $+0.068$ macro-F1, with FedProx
winning all $10$ paired seeds. This reproduces the qualitative
finding of Li et al.\ (2020), §5.2 on the DermaMNIST task: the
canonical asymmetric straggler setup produces a substantially
higher macro-F1 under FedProx than under FedAvg. The magnitude
is approximately four times larger than the symmetric C2
$\Delta$ of $+0.017$ reported in
Section~\ref{sec:sh-engineered-results}, primarily because the
asymmetric setup makes two simultaneous changes to FedAvg
relative to FedProx (dropping partial updates and lacking the
proximal regulariser) rather than only one.

% ==============================================================================
\section{Partial Client Participation}
\label{sec:sh-partial-participation}
% ==============================================================================

The system-heterogeneity conditions reported in
Sections~\ref{sec:sh-engineered-results}--\ref{sec:sh-asymmetric}
vary the second canonical FL parameter of system heterogeneity ---
the local-epoch budget $E$ --- while holding the per-round
participation fraction fixed at $C = 1.0$ (every client
participates every round). The third canonical FL parameter, $C$,
is itself an axis of system heterogeneity widely tested in the
literature. McMahan et al.\ (2017) vary
$C \in \{0.1, 0.2, 0.5, 1.0\}$ in the original FedAvg paper;
Li et al.\ (2020), §5.3 test FedProx under partial participation
on Synthetic and FEMNIST; Kairouz et al.\ (2021) list $C$ among
the defining FL setup parameters. In real federated deployments,
not all institutions respond every round, so the full-participation
setting used so far in this chapter is the cleanest mathematical
baseline but the least operationally realistic case. This section
reports the FedAvg-vs-FedProx comparison under partial client
participation at $C = 0.5$ (four of seven clients sampled
uniformly at random per round), holding all other hyperparameters
identical to the C0 engineered baseline of
Section~\ref{sec:sh-engineered-results}.

\paragraph{Mechanism prediction.}
Under partial participation, the per-round aggregation pool is
smaller and its composition varies across rounds. Any single
client's drift therefore contributes a larger fraction of the
aggregated update than under full participation, and the
round-to-round aggregation variance is higher. The proximal
anchor's drift-control mechanism is theoretically more useful in
this regime: with more drift to constrain and more aggregation
variance to smooth, the proximal regulariser is expected to
exert a stronger relative restraint on per-client updates than at
$C = 1.0$. The expected direction is positive; the magnitude is
not strongly predicted by the theory and is reported
descriptively below.

\subsection{Setup}
\label{sec:sh-partial-setup}

At each communication round, $\lfloor 0.5 \times 7 \rfloor + 1 = 4$
clients are sampled uniformly at random (without replacement)
from the seven available clients; the remaining three clients do
not participate in that round and their parameters are not
updated. All other hyperparameters are held identical to the
engineered C0 baseline (Section~\ref{sec:sh-engineered-results}):
the same \texttt{balanced\_paired\_7\_clients} partition, $E = 20$
local epochs, $R = 150$ communication rounds, FedProx proximal
coefficient $\mu = 0.01$, SGD with learning rate $0.01$ and
momentum $0.9$, batch size $32$, and the same ten random seeds.
The Flower simulation runtime is used.

\paragraph{A note on the per-round sampling realisation.}
The two algorithm arms share the same starting random seed, and
on $5$ of the $10$ seeds the per-round sampled client set is
identical across the FedAvg and FedProx arms for all $150$
rounds. On the remaining $5$ seeds the sampling streams desync
from a small number of rounds onward, due to FedProx's code path
consuming one additional random number per round. We therefore
do not rely on a strict paired-comparison interpretation of the
two arms in this section. Outcome-level (test macro-F1) results
are reported as two-arm summaries (mean $\pm$ standard deviation
across the ten seeds, per algorithm), and the load-bearing
mechanism claim of this section rests on the per-client
update-norm comparison
(Section~\ref{sec:sh-partial-update-norms}), which is a within-
arm scalar and is robust to differences in which specific
clients participated in any given round.

\subsection{Results}
\label{sec:sh-partial-results}

Table~\ref{tab:sh-partial} reports the test macro-F1 at the
best-validation checkpoint under partial participation, side by
side with the full-participation engineered C0 baseline from
Section~\ref{sec:sh-engineered-results}.
Figure~\ref{fig:sh-partial-convergence} shows the corresponding
per-round validation macro-F1 trajectories.

\begin{table}[!htbp]
\centering
\caption{Partial-participation result on the engineered partition
($n = 10$ seeds per arm, Flower runtime). The first row
reproduces the full-participation baseline from
Section~\ref{sec:sh-engineered-results}; the second row reports
the partial-participation result. Test macro-F1 is reported as
mean $\pm$ standard deviation across seeds per arm. The
$\Delta$ column is the difference of arm means, reported
descriptively (see the setup note for why a paired-comparison
interpretation is not claimed here). The higher arm mean per
row is shown in bold.}
\label{tab:sh-partial}
\small
\begin{tabular}{lccc}
\toprule
Setting & FedAvg & FedProx & $\Delta$ (mean) \\
\midrule
$C = 1.0$ (full participation)         & $0.4964 \pm 0.0150$ & $\mathbf{0.5033 \pm 0.0206}$ & $+0.007$ \\
$C = 0.5$ (partial, 4/7 per round)     & $0.4780 \pm 0.0234$ & $\mathbf{0.4944 \pm 0.0185}$ & $+0.016$ \\
\bottomrule
\end{tabular}
\end{table}

Mean test macro-F1 is higher for FedProx than FedAvg under both
participation regimes, with the arm gap approximately doubling
from $+0.007$ at $C = 1.0$ to $+0.016$ at $C = 0.5$. The direction
is consistent with the mechanism prediction: with the per-round
aggregation pool reduced from seven clients to four, the proximal
anchor's drift-control role becomes more useful at the
parameter-level (see Section~\ref{sec:sh-partial-update-norms})
and the outcome-level arm gap enlarges in parallel. The
across-seed standard deviations on both arms ($0.02$--$0.03$
macro-F1) are of the same order as the arm-gap magnitudes,
so the macro-F1 numbers are reported here as supporting
descriptive evidence; the load-bearing mechanism claim is the
update-norm result reported in
Section~\ref{sec:sh-partial-update-norms}.

\begin{figure}[!htbp]
\centering
\includegraphics[width=\textwidth]{F_partial_participation.pdf}
\caption{Validation macro-F1 over $150$ communication rounds on
the engineered partition under full and partial client
participation (Flower runtime, $n = 10$ seeds per arm, mean
$\pm$ SEM band). \textbf{Left ($C = 1.0$):} all seven clients
participate every round; FedProx maintains a modest separation
from FedAvg. \textbf{Right ($C = 0.5$):} four of seven clients
are sampled uniformly at random per round; the gap between the
two arm means is visibly wider than under full participation,
with FedProx's arm mean ($0.494$ test macro-F1) above FedAvg's
($0.478$) throughout late training.}
\label{fig:sh-partial-convergence}
\end{figure}

\paragraph{Algorithmic robustness to partial participation.}
A complementary observation is that the two algorithms differ in
how much absolute test macro-F1 they lose when participation is
reduced from $C = 1.0$ to $C = 0.5$. FedAvg loses $0.018$ macro-F1
in absolute terms ($0.496 \to 0.478$), while FedProx loses only
$0.009$ ($0.503 \to 0.494$). The proximal anchor therefore
makes FedProx approximately twice as robust to the reduction in
participation fraction as FedAvg in this setup. Both algorithms
remain operational under partial participation; the difference
is the rate at which performance degrades as the participation
fraction shrinks.

\subsection{Per-class breakdown}
\label{sec:sh-partial-per-class}

Table~\ref{tab:sh-partial-per-class} decomposes the
partial-participation arm comparison by class. The per-class
pattern matches the structure observed under both the engineered
C0 baseline and the C2 random-straggler condition: melanoma is
the dominant contributor to the macro-F1 arm gap, and the
universal class melanocytic nevi is mechanistically inert.

\begin{table}[!htbp]
\centering
\caption{Per-class test F1 under partial participation $C = 0.5$
on the engineered partition (Flower runtime, $n = 10$ seeds per
arm, mean $\pm$ standard deviation across seeds per arm). The
higher arm mean per row is shown in bold; $\Delta$ is the
difference of arm means (FedProx $-$ FedAvg) reported
descriptively.}
\label{tab:sh-partial-per-class}
\small
\begin{tabular}{lccc}
\toprule
Class (prevalence) & FedAvg & FedProx & $\Delta$ (mean) \\
\midrule
Actinic keratoses (3.27\%)        & $0.453 \pm 0.040$ & $\mathbf{0.462 \pm 0.044}$ & $+0.008$ \\
Basal cell carcinoma (5.13\%)     & $0.454 \pm 0.045$ & $\mathbf{0.472 \pm 0.040}$ & $+0.018$ \\
Benign keratosis (10.97\%)        & $0.403 \pm 0.069$ & $\mathbf{0.410 \pm 0.052}$ & $+0.006$ \\
Dermatofibroma (1.15\%)           & $0.238 \pm 0.093$ & $\mathbf{0.277 \pm 0.077}$ & $+0.039$ \\
Melanoma (11.11\%)                & $0.199 \pm 0.063$ & $\mathbf{0.282 \pm 0.033}$ & $\mathbf{+0.084}$ \\
Melanocytic nevi (67.05\%)        & $\mathbf{0.881 \pm 0.007}$ & $0.882 \pm 0.004$ & $+0.002$ \\
Vascular lesions (1.41\%)         & $\mathbf{0.718 \pm 0.036}$ & $0.675 \pm 0.052$ & $-0.043$ \\
\bottomrule
\end{tabular}
\end{table}

The melanoma per-class arm gap under partial participation
($\Delta = +0.084$) is almost identical to that observed under
the C2 random-straggler condition
($\Delta = +0.085$; Section~\ref{sec:sh-c2-per-class}) and larger
than the C0 full-participation baseline
($\Delta = +0.068$; Section~\ref{sec:engineered-per-class}). The
melanocytic-nevi $\Delta$ remains effectively zero
($+0.002$, both arms within $\pm 0.007$ across seeds), preserving
the mechanism-null pattern observed in every previous condition.
The clinical localisation of the FedProx advantage on the
melanoma class therefore generalises across both axes of system
heterogeneity (variable local epochs and variable client
participation).

The single class on which FedProx's arm mean is lower than
FedAvg's under partial participation is \emph{vascular lesions}
($\Delta = -0.043$). Vascular lesions are held by clients C4 and
C5 in the engineered partition (the melanoma/vascular specialty
pair). Under $C = 0.5$, both C4 and C5 are sampled out of any
given round with joint probability
$\binom{5}{4} / \binom{7}{4} \approx 14\%$, so vascular receives
no training signal in roughly one in every seven rounds. The
small absolute count of vascular training samples
($n = 99$ globally) means the model has limited capacity to
recover from these zero-signal rounds. This class-specific
regression is a deployment-relevant observation: in production
federations with rare classes concentrated on a small number of
clients, partial participation may degrade performance on those
classes even where the macro-F1 aggregate improves.

\subsection{Mechanism evidence: update norms}
\label{sec:sh-partial-update-norms}

The parameter-level mechanism evidence introduced in
Section~\ref{sec:update-norms} can also be measured under
partial participation. This is the primary inferential result
of this section: the update-norm comparison is a within-arm
scalar (mean drift size per algorithm) and does not depend on
which specific clients participated in any individual round, so
it is unaffected by the per-round sampling realisation
described in the setup. Table~\ref{tab:sh-partial-norms} reports
the mean per-client update norm $\lVert w_i^{t+1} - w^t \rVert_2$
under each participation regime.

\begin{table}[!htbp]
\centering
\caption{Per-client update norm averaged across all (round,
client) pairs and across paired seeds, on the engineered
partition (Flower runtime). The ratio column reports the FedAvg
mean update norm divided by the FedProx mean update norm under
each participation regime.}
\label{tab:sh-partial-norms}
\small
\begin{tabular}{lccc}
\toprule
Setting & FedAvg $\lVert \Delta w \rVert$ & FedProx $\lVert \Delta w \rVert$ & Ratio FA/FP \\
\midrule
$C = 1.0$ (full participation)    & $1.523$ & $\mathbf{1.029}$ & $1.48\times$ \\
$C = 0.5$ (partial participation) & $1.669$ & $\mathbf{1.078}$ & $\mathbf{1.55\times}$ \\
\bottomrule
\end{tabular}
\end{table}

Two observations follow.

First, both algorithms produce larger per-client update norms
under partial participation than under full participation
(FedAvg $1.52 \to 1.67$, FedProx $1.03 \to 1.08$). With fewer
clients aggregating per round, each remaining client's local
trajectory must travel further to ``absorb'' the global model's
adjustment toward its local distribution, producing larger
per-round parameter updates on both algorithms.

Second, the FedAvg-to-FedProx update-norm ratio \emph{increases}
from $1.48\times$ under full participation to $1.55\times$ under
partial participation. The proximal regulariser is empirically
more active in restraining client drift when the per-round
aggregation pool is smaller --- the regime in which the
mechanism prediction said it should be most useful. The
parameter-level signature of the mechanism therefore strengthens
in the same direction as the outcome-level $\Delta$ enlarges:
both signals move together, in the direction predicted before
running the experiment.

\subsection{Summary}
\label{sec:sh-partial-summary}

Under partial client participation at $C = 0.5$ on the
engineered partition, the parameter-level mechanism evidence
strengthens relative to the full-participation baseline: the
FedAvg-to-FedProx mean update-norm ratio enlarges from
$1.48\times$ to $1.55\times$ --- the predicted direction, given
that a smaller per-round aggregation pool leaves more client
drift for the proximal anchor to constrain. This is the primary
finding of the section and is robust to the per-round sampling
realisation, since the update-norm ratio compares within-arm
average drift sizes rather than within-pair outcome differences.
The outcome-level (test macro-F1) results move in parallel: the
mean arm gap approximately doubles ($+0.007 \to +0.016$
macro-F1), with FedProx's arm mean higher than FedAvg's under
both regimes. The per-class structure is preserved across
regimes: melanoma is the dominant contributor to the arm gap
($\Delta = +0.084$ at $C = 0.5$, comparable to $+0.085$ under
C2 random stragglers and larger than $+0.068$ at $C = 1.0$),
and the universal class melanocytic nevi remains mechanism-inert
($\Delta = +0.002$). The single per-class regression
(vascular lesions, $\Delta = -0.043$) is mechanism-coherent with
the engineered partition's concentration of vascular samples on
the specialty pair (C4, C5) and is reported as a
deployment-relevant caveat rather than a contradiction.

% ==============================================================================
% Future sections to be added below in subsequent writing iterations:
%
%   \section{Statistical Heterogeneity}
%     \subsection{Experimental design and partition motivation}
%     \subsection{IID baseline (mechanism null)}
%     \subsection{Engineered partition (primary results)}
%     \subsection{Per-class breakdown}
%     \subsection{Cross-runtime sensitivity}
%     \subsection{Robustness across partition designs}
%     \subsection{Direct mechanism evidence: update norms}
%     \subsection{Federation tax}
%
%   \section{System Heterogeneity}
%     \subsection{Experimental design and condition motivation}
%     \subsection{Results on the engineered partition}
%     \subsection{Cross-partition control: IID partition}
%     \subsection{Partial client participation}
%     \subsection{Replication of the canonical Li et al. 5.2 protocol}
%
%   \section{Summary and limitations}
% ==============================================================================

\end{document}
